\chapter{Interpretable Regression with Checked Physical Predictions}
\label{ch:interpretable}

\section{Rationale}

A component prediction is useful for engineering analysis when its dependence on the inputs can be examined and its material accounting can be checked. These requirements address different questions. A transparent regression equation explains how the predicted state depends on operating conditions and influent loads. A mass conservation test determines whether that state preserves the inventories represented by the reaction network. A non-negativity test determines whether each predicted concentration is physically meaningful. An explicit equation can fail either physical test. A physically projected prediction can also remain difficult to interpret if the underlying regression is opaque.

Invariant-constrained second-order regression (ICSOR) combines an explicit polynomial driver, learned dependence among effluent components, and checked prediction. A polynomial driver is the fitted expression containing inputs, squares, and pairwise products before component coupling is applied. In this chapter, deployment means evaluating the fitted model at a new influent and operating setting and returning a physically checked component state. The driver retains separate terms for operating conditions, influent composition, curvature, and interactions. The coupling relation expresses how the driver contributes to the predicted component state. Deployment then checks that state and applies a constrained projection when necessary. Together, these stages provide an inspectable regression and a checked effluent state.

The preference for an inspectable regression follows the argument of \citet{Rudin2019} that transparent model structures can support decisions that require scrutiny. It also follows the broader development of models that combine learned relations with physical structure discussed by \citet{Shen2023}. The fitted relations are available for examination, comparison with process knowledge, and testing across operating conditions.

Activated sludge nutrient removal provides a strong reason to retain interactions. Carbon availability influences denitrification and the activity of organisms that store phosphorus. Oxygen availability influences both carbon oxidation and nitrification. \citet{Wentzel1992} describe these connected pathways, and \citet{Guerrero2011} examine the role of carbon source in the competition between phosphorus-storing organisms and denitrifiers. A regression that represents an aeration effect independently of loading can overlook such dependencies. An interaction term can describe an association between an operating action and a loading condition while keeping that association visible.

The prediction target is the effluent component vector \(c_{out}\in\mathbb{R}^{F}\). The adopted process basis has \(F=20\) components and \(P=28\) reactions. The operating vector \(u\in\mathbb{R}^{M_{op}}\) has \(M_{op}=2\) entries for HRT and aeration level. The full influent vector is \(c_{in}\in\mathbb{R}^{F}\). Component concentrations retain their physical units. The four reported composites are obtained from a fixed map \(I_{\mathrm{comp}}\in\mathbb{R}^{4\times F}\). That map is applied after component prediction.

\section{Theory}

\subsection{Mass conservation in the adopted state space}

The link between a component balance and a regression constraint can first be seen in a hypothetical conversion of one component into another. Suppose that one unit of component 1 is converted into one unit of component 2. The two components are expressed in a common inventory unit for this illustration. Their changes satisfy
\begin{equation}
\equationname{Illustrative component changes for a two-component reaction}
\label{eq:icsor_scalar_inventory_example}
c_{out,1}-c_{in,1}=-\xi,
\qquad
c_{out,2}-c_{in,2}=\xi.
\end{equation}
Adding these equations eliminates the unknown conversion amount. The relation available to the regression is \(c_{out,1}+c_{out,2}=c_{in,1}+c_{in,2}\). With an illustrative influent \((6,4)^\top\) and conversion amount 2, the effluent is \((4,6)^\top\). Both totals are 10. Many other non-negative pairs also have that total. The balance defines a set of regression outputs with the same total inventory.

The same calculation can be written using the small matrices
\begin{equation}
\equationname{Illustrative stoichiometric matrix and mass conservation operator}
\label{eq:icsor_small_inventory_matrices}
\nu=\begin{bmatrix}-1&1\end{bmatrix},
\qquad
A=\begin{bmatrix}1&1\end{bmatrix},
\qquad
A\nu^\top=1(-1)+1(1)=0.
\end{equation}
These hypothetical matrices use reduced dimensions to make each operation visible. The adopted reactor uses the stated 20-component basis, reaction matrix, and operating domain.

The stoichiometric matrix \(\nu\in\mathbb{R}^{P\times F}\) records reactions as rows and components as columns. Let \(V_0\in\mathbb{R}^{F\times K}\) span its right null space, where \(K=5\). The invariant operator is
\begin{equation}
\equationname{Structured-regression mass conservation operator from the null space}
\label{eq:icsor_invariant_operator}
A=V_0^\top\in\mathbb{R}^{K\times F},
\qquad \nu V_0=0,
\qquad A\nu^\top=0.
\end{equation}
The columns of \(V_0\) are independent, so \(A\) has full row rank. With an orthonormal null-space basis, \(AA^\top=I_K\). A different full-rank basis describes the same mass conservation conditions, although its row scaling changes the numerical size of a soft invariant penalty.

For the control volume and boundary accounting represented by the adopted process model, a reaction-induced change has the form
\begin{equation}
\equationname{Effluent change in stoichiometric process coordinates}
\label{eq:icsor_stoichiometric_change}
c_{out}-c_{in}=\nu^\top\xi,
\qquad \xi\in\mathbb{R}^{P}.
\end{equation}
The reaction-progress vector \(\xi\) expresses the net contribution of the reactions in concentration-equivalent units. Multiplication by \(A\) removes that unknown contribution and gives
\begin{equation}
\equationname{Mass conservation equality for the structured-regression effluent}
\label{eq:icsor_invariant_balance}
A(c_{out}-c_{in})=0.
\end{equation}
This construction uses the mass conservation logic of the photochemical surrogate framework of \citet{SturmWexler2020}. \citet{SturmWexler2022} further show how a stoichiometric transformation can connect learned reaction information to state changes satisfying mass conservation. The same linear algebra applies here, but the component combinations used for mass conservation are those of the activated sludge reaction network.

Equation~\eqref{eq:icsor_invariant_balance} uses the modeled boundary in Equation~\eqref{eq:icsor_stoichiometric_change}. External gas transfer, chemical dosing, solids withdrawal, and other boundary fluxes enter the balance where they occur. With an explicit boundary contribution \(s\), the relation is \(A(c_{out}-c_{in}-s)=0\).

For a new influent state, the component set used by checked deployment is
\begin{equation}
\equationname{Structured-regression feasible set for mass conservation and non-negativity}
\label{eq:icsor_feasible_set}
\mathcal{F}(c_{in})=
\left\{c\in\mathbb{R}^{F}\mid Ac=Ac_{in},\ c\geq\mathbf{0}\right\}.
\end{equation}
When \(c_{in}\geq\mathbf{0}\), this set contains \(c_{in}\) and is therefore nonempty. Its members satisfy the modeled linear invariants and component non-negativity within the declared boundary.

The reported composites provide another reason to work in component space. Distinct component states can have the same COD, TN, TP, and TSS values. A negative component can also be hidden by positive contributions to an aggregate. When \(I_{\mathrm{comp}}\) is entrywise non-negative, a non-negative component state produces non-negative composites. The reverse implication does not generally hold. Mass conservation and non-negativity are therefore tested before the component state is collapsed into the four reported quantities.

\subsection{A partitioned second-order driver}

\subsubsection{Separating operation from loading}

Operating conditions and influent composition play different engineering roles. HRT and aeration define aspects of the reactor environment. Influent components define the material presented to that environment. This distinction is consistent with the operating-sensitivity analysis of \citet{Araujo2013} and the separation of operational and influent uncertainty examined by \citet{Machado2014}.

Start with two hypothetical operating entries \(u_1,u_2\) and two influent entries \(c_1,c_2\). A second-order model uses the individual entries, their squares, and their pairwise products. The operating products are \(u_1^2,u_1u_2,u_2u_1,u_2^2\). The feature vector keeps both orders as an accounting convention, with \(u_1u_2=u_2u_1\).

For two column vectors, the Kronecker operation multiplies the entire second vector by the first entry of the first vector, then by its second entry. Thus
\begin{equation}
\equationname{Illustrative operating, influent, and mixed Kronecker products}
\label{eq:icsor_small_kronecker}
\begin{aligned}
u\otimes u&=\begin{bmatrix}u_1u_1\\u_1u_2\\u_2u_1\\u_2u_2\end{bmatrix},
&c\otimes c&=\begin{bmatrix}c_1c_1\\c_1c_2\\c_2c_1\\c_2c_2\end{bmatrix},\\
u\otimes c&=\begin{bmatrix}u_1c_1\\u_1c_2\\u_2c_1\\u_2c_2\end{bmatrix}.
\end{aligned}
\end{equation}
For the illustrative values \(u=(2,3)^\top\) and \(c=(4,5)^\top\), these three blocks are respectively \((4,6,6,9)^\top\), \((16,20,20,25)^\top\), and \((8,10,12,15)^\top\). Including the baseline 1 and the four first-order entries gives 17 ordered features in this reduced example. The same construction applies to the larger adopted input vector.

The feature vector retains both roles through the ordered expansion
\begin{equation}
\equationname{Second-order driver feature vector and its dimension}
\label{eq:icsor_features}
\phi(u,c_{in})=
\begin{bmatrix}
1\\u\\c_{in}\\u\otimes u\\c_{in}\otimes c_{in}\\u\otimes c_{in}
\end{bmatrix}
\in\mathbb{R}^{D},
\qquad
D=1+M_{op}+F+M_{op}^{2}+F^{2}+M_{op}F.
\end{equation}
The symbol \(\otimes\) denotes the Kronecker product. The adopted dimensions give \(D=467\). The ordered expansion includes both \(u_i u_j\) and \(u_j u_i\), and both \(c_{in,i}c_{in,j}\) and \(c_{in,j}c_{in,i}\). These pairs are equal as numerical features. Their interpretation requires the symmetry convention developed below.

Multiplication by regression coefficients can also be read one entry at a time. For an illustrative first driver, take the baseline coefficient as 1, operating coefficients as \((2,-1)\), and influent coefficients as \((0.5,0)\). Let its operating quadratic row be \((0.1,0.1,0.1,0)\), influent quadratic row be \((0,0,0,-0.05)\), and operating-load row be \((0,0,0.25,0)\). Applying these rows to the numerical feature blocks gives
\begin{equation}
\equationname{Numerical evaluation of an illustrative regression driver}
\label{eq:icsor_driver_multiplication_example}
\begin{aligned}
r_1={}&1+\begin{bmatrix}2&-1\end{bmatrix}\begin{bmatrix}2\\3\end{bmatrix}
+\begin{bmatrix}0.5&0\end{bmatrix}\begin{bmatrix}4\\5\end{bmatrix}\\
&+\begin{bmatrix}0.1&0.1&0.1&0\end{bmatrix}\begin{bmatrix}4\\6\\6\\9\end{bmatrix}\\
&+\begin{bmatrix}0&0&0&-0.05\end{bmatrix}\begin{bmatrix}16\\20\\20\\25\end{bmatrix}\\
&+\begin{bmatrix}0&0&0.25&0\end{bmatrix}\begin{bmatrix}8\\10\\12\\15\end{bmatrix}\\
={}&1+1+2+1.6-1.25+3=7.35.
\end{aligned}
\end{equation}
Each row-vector multiplication is a sum of coefficient times feature. It is the same operation as an ordinary linear regression, but the supplied features include nonlinear products. A second hypothetical driver \(r_2=-u_1+u_2+0.2c_2\) would equal 2 at the same inputs. The coupled relation converts the two driver values into predicted concentrations.

Before grouping the coefficients into matrices, the first driver in this illustration can be written as
\begin{equation}
\equationname{Scalar polynomial form of the illustrative regression driver}
\label{eq:icsor_scalar_driver_example}
r_1=1+2u_1-u_2+0.5c_1+0.1u_1^2
+0.2u_1u_2-0.05c_2^2+0.25u_2c_1.
\end{equation}
The coefficient 0.2 of \(u_1u_2\) is the sum of the two ordered coefficients 0.1. In a general row, the scalar calculation is \(r_f=\sum_{\ell=1}^{D}B_{f\ell}\phi_\ell\). Carrying out this sum for all components gives the driver \(r\in\mathbb{R}^{F}\) through
\begin{equation}
\equationname{Matrix mapping from second-order features to the regression driver}
\label{eq:icsor_driver}
r(u,c_{in})=B\phi(u,c_{in}),
\qquad B\in\mathbb{R}^{F\times D}.
\end{equation}
The driver is an intermediate regression quantity supplied to the coupled prediction relation. Its coefficients have no sign restriction, so the regression can represent both increasing and decreasing effects. Non-negativity is imposed on component states where it has a physical meaning. The driver block form is
\begin{equation}
\equationname{Driver decomposition into operating, influent, and interaction blocks}
\label{eq:icsor_driver_blocks}
\begin{aligned}
r(u,c_{in})={}&b+W_u u+W_{in}c_{in}
+\Theta_{uu}(u\otimes u)\\
&+\Theta_{cc}(c_{in}\otimes c_{in})
+\Theta_{uc}(u\otimes c_{in}).
\end{aligned}
\end{equation}
Here \(b\in\mathbb{R}^{F}\) is a baseline term. The matrices \(W_u\in\mathbb{R}^{F\times M_{op}}\) and \(W_{in}\in\mathbb{R}^{F\times F}\) contain first-order operating and influent terms. The matrices \(\Theta_{uu}\in\mathbb{R}^{F\times M_{op}^{2}}\) and \(\Theta_{cc}\in\mathbb{R}^{F\times F^{2}}\) contain curvature within each input group. The matrix \(\Theta_{uc}\in\mathbb{R}^{F\times M_{op}F}\) contains operating-load interactions.

\subsubsection{Illustrative operating and loading effects}

Consider an illustrative driver for one component with aeration level \(a\) and influent ammonium concentration \(n\). A term of the form
\begin{equation}
\equationname{Illustrative driver with aeration, nitrogen, and interaction effects}
\label{eq:icsor_example_driver}
r_f=\beta_0+\beta_a a+\beta_n n+\beta_{aa}a^2+\beta_{an}an
\end{equation}
has an aeration derivative \(\partial r_f/\partial a=\beta_a+2\beta_{aa}a+\beta_{an}n\). The effect described by this derivative changes with both aeration level and ammonium loading. The derivative of this hypothetical polynomial combines the linear aeration term, curvature, and the aeration-ammonium interaction.

The same reasoning applies to simultaneous influent loads. A carbon-ammonium interaction allows the modeled association with carbon availability to depend on nitrogen loading. A phosphorus-storage interaction allows the modeled association with dissolved phosphate to depend on a particulate storage pool. These terms provide a transparent approximation to a nonlinear response involving the microbial pathways described by \citet{Henze1999}.

The second-order library represents smooth local curvature and pairwise input interactions. Explicit nonlinear libraries make candidate response structures available for inspection, as illustrated by \citet{Brunton2016}.

\subsubsection{Symmetry and the meaning of cross terms}

For a two-entry vector, direct multiplication gives
\begin{equation}
\equationname{Expansion of a two-variable quadratic form}
\label{eq:icsor_scalar_quadratic}
\begin{bmatrix}x_1&x_2\end{bmatrix}
\begin{bmatrix}q_{11}&q_{12}\\q_{21}&q_{22}\end{bmatrix}
\begin{bmatrix}x_1\\x_2\end{bmatrix}
=q_{11}x_1^2+(q_{12}+q_{21})x_1x_2+q_{22}x_2^2.
\end{equation}
Only the sum of the two off-diagonal coefficients enters this polynomial. Replacing them by their average leaves that sum unchanged. For example,
\begin{equation}
\equationname{Numerical example of quadratic-coefficient symmetrization}
\label{eq:icsor_symmetry_numbers}
Q=\begin{bmatrix}0.1&0.3\\-0.1&0\end{bmatrix},
\qquad
\frac{Q+Q^\top}{2}=\begin{bmatrix}0.1&0.1\\0.1&0\end{bmatrix}.
\end{equation}
At \(x=(2,3)^\top\), the first matrix gives \(0.4+1.8-0.6=1.6\). The symmetric matrix gives \(0.4+0.6+0.6=1.6\). The asymmetric coefficients cannot be distinguished from these predictions. Their squared sums are different. The original matrix has squared sum \(0.01+0.09+0.01=0.11\), while its symmetric part has squared sum \(0.01+0.01+0.01=0.03\). The removed part has squared sum 0.08 and contributes no prediction. This is why a positive coefficient penalty favors the symmetric coefficient matrix.

For component \(f\), reshape the corresponding rows of \(\Theta_{uu}\) and \(\Theta_{cc}\) into square matrices \(Q_{uu}^{(f)}\) and \(Q_{cc}^{(f)}\). A quadratic term can then be written as \(x^\top Qx\). For any real square matrix,
\begin{equation}
\equationname{Symmetric coefficient matrix of a quadratic form}
\label{eq:icsor_quadratic_symmetry}
x^\top Qx=x^\top\operatorname{sym}(Q)x,
\qquad
\operatorname{sym}(Q)=\frac{Q+Q^\top}{2}.
\end{equation}
The skew-symmetric part contributes nothing to the prediction. It can nevertheless increase the coefficient norm because
\begin{equation}
\equationname{Frobenius-norm decomposition into symmetric and skew parts}
\label{eq:icsor_quadratic_norm}
\|Q\|_F^2=
\|\operatorname{sym}(Q)\|_F^2+
\|\operatorname{skew}(Q)\|_F^2.
\end{equation}
Replacing each square slice by its symmetric part preserves every fitted value and cannot increase a positive ridge penalty. The projection is applied after each driver update. With a strictly positive driver ridge weight, the exact minimum-norm solution has no skew-symmetric part.

Symmetry also determines how a coefficient display is read. If \(q_{ij}=q_{ji}\), the two displayed off-diagonal cells together contribute \(2q_{ij}x_i x_j\). A displayed cell is therefore half the coefficient of the corresponding unordered cross term. For an illustrative matrix with \(q_{12}=q_{21}=0.3\), the polynomial contribution is \(0.6x_1x_2\). Reading either cell as the entire interaction coefficient would halve its contribution.

Each component has \(M_{op}(M_{op}+1)/2=3\) distinct operating quadratic terms and \(F(F+1)/2=210\) distinct influent quadratic terms. Together with the baseline, linear terms, and operating-load interactions, the library has 276 distinct polynomial terms per component. These counts define the structural dimensions of the polynomial library.

\subsection{Learned coupling among effluent components}

Consider two hypothetical component equations in which each channel has a coupling coefficient 0.2 to the other channel. With driver values 4 and 2, their equations are
\begin{equation}
\equationname{Illustrative coupled equations for two effluent components}
\label{eq:icsor_scalar_coupling_example}
c_1=4+0.2c_2,
\qquad c_2=2+0.2c_1.
\end{equation}
Substituting the second equation into the first gives \(c_1=4+0.2(2+0.2c_1)\). Moving the remaining \(c_1\) term to the left gives \(0.96c_1=4.4\). Therefore \(c_1=55/12\) and \(c_2=35/12\). The solution is approximately \((4.5833,2.9167)^\top\). These hypothetical values demonstrate the arithmetic of a simultaneous regression relation.

Moving the coupled terms to the left makes the same calculation a matrix solve,
\begin{equation}
\equationname{Illustrative coupling matrix and coupled-state system}
\label{eq:icsor_small_coupling_matrices}
\Gamma=\begin{bmatrix}0&0.2\\0.2&0\end{bmatrix},
\qquad
R=\begin{bmatrix}1&-0.2\\-0.2&1\end{bmatrix},
\qquad
R\begin{bmatrix}c_1\\c_2\end{bmatrix}=\begin{bmatrix}4\\2\end{bmatrix}.
\end{equation}
The determinant of \(R\) is \(1-0.2^2=0.96\). Its inverse is
\begin{equation}
\equationname{Inverse of the illustrative two-component coupling system}
\label{eq:icsor_small_coupling_inverse}
R^{-1}=\frac{1}{0.96}\begin{bmatrix}1&0.2\\0.2&1\end{bmatrix}.
\end{equation}
Multiplying this inverse by the driver gives the same two scalar solutions. The inverse also explains why a change in one driver can change both predicted components.

The regression includes a matrix \(\Gamma\in\mathbb{R}^{F\times F}\) that describes fitted dependence among the effluent components. Define
\begin{equation}
\equationname{General coupled-state relation and system matrix}
\label{eq:icsor_coupled_relation}
R=I_F-\Gamma,
\qquad Rc=r(u,c_{in}).
\end{equation}
Equivalently, the equation for component \(f\) is \(c_f=r_f+\sum_{j\neq f}\Gamma_{fj}c_j\). This form separates the input-dependent driver from dependence among the predicted component channels. The diagonal of \(\Gamma\) is fixed to zero to avoid a redundant self-coupling coefficient.

The sign of \(\Gamma_{fj}\) describes reinforcement or compensation within the fitted simultaneous relation. Activated sludge reactions involve several components at once, and the stoichiometric matrix specifies those transformations. The coupling matrix is learned from state patterns. Its entries describe fitted dependence associated with correlated transformations, loading patterns, and the allocation of effects between the driver and coupling blocks.

When \(R\) is nonsingular, the raw component prediction is \(R^{-1}r\). Multiplication by \(R^{-1}\) redistributes or amplifies the response represented by the driver. For interpretation, the driver, coupling matrix, and complete raw mapping are examined together.

The adopted coupling constraints require
\begin{equation}
\equationname{Diagonal, magnitude, and conditioning safeguards for component coupling}
\label{eq:icsor_coupling_safeguards}
\Gamma_{ff}=0,
\qquad |\Gamma_{fj}|\leq\gamma_{abs}\quad(f\neq j),
\qquad \kappa_2(I_F-\Gamma)\leq\kappa_{max}.
\end{equation}
The condition number is the ratio of the largest to smallest singular value of the coupled-system matrix. A nearly singular matrix can amplify small driver perturbations and numerical errors. The condition check controls the complete matrix separately from the bounds on individual entries.

The small symmetric matrix in Equation~\eqref{eq:icsor_small_coupling_matrices} illustrates this measure without a large singular-value calculation. In the direction \((1,1)^\top\), multiplication by \(R\) scales the vector by 0.8. In the direction \((1,-1)^\top\), it scales the vector by 1.2. These orthogonal directions give singular values 0.8 and 1.2, so the condition number is \(1.2/0.8=1.5\).

If the two illustrative couplings were 0.99, the same directions would have scales 0.01 and 1.99. The condition number would be 199. A driver error of magnitude \(\epsilon\) in each component along the first direction would produce a state error of \(100\epsilon\) in each component. This hypothetical choice illustrates amplification at a coupling magnitude of 0.99, above the adopted bound.

Even a small entry bound can permit singularity in a larger matrix. In a hypothetical four-component matrix with every off-diagonal entry equal to \(1/3\), every row sum of \(\Gamma\) is one. Hence \(\Gamma\mathbf{1}=\mathbf{1}\) and \((I-\Gamma)\mathbf{1}=0\). Each entry is smaller than the adopted bound 0.368, but the coupled-system matrix is singular. This example explains why entry bounds and the condition-number test perform separate tasks.

Shrinking a candidate coupling matrix toward zero can provide a safe candidate because \(I_F-0=I_F\). A zero coupling matrix gives a condition number of one. A prescribed \(\kappa_{max}>1\) therefore permits this fallback. The nonlinear condition-number requirement defines a generally nonconvex admissible coupling set. It is evaluated after the convex box-constrained estimation problem proposes a coupling update.

\section{Methodology}

\subsection{Fitting in prediction space}

\subsubsection{The coupled training objective}

The fitting problem adjusts coefficients to approximate the mechanistic effluent concentrations. It also uses auxiliary fitted states, which are concentration variables optimized for the training cases. These variables help express fitting constraints and penalties. They differ from the component predictions calculated for a new influent after fitting.

The squared error of one prediction is the square of its difference from the reference. For several components and observations, fitting adds these squared differences. A compact matrix norm represents this sum. For any matrix \(M\) with entries \(M_{ij}\),
\begin{equation}
\equationname{Squared Frobenius norm as a sum of coefficient squares}
\label{eq:icsor_frobenius_definition}
\|M\|_F^2=\sum_i\sum_j M_{ij}^2.
\end{equation}
For example, a residual matrix \(\begin{bmatrix}1&-2\\3&0\end{bmatrix}\) has squared norm \(1^2+(-2)^2+3^2+0^2=14\). The norm is \(\sqrt{14}\), while a least-squares objective uses 14. Squaring prevents positive and negative errors from canceling when they are added.

Let \(i\) identify an observation, \(f\) a component, \(k\) an invariant, and \(\ell\) a feature. The entry \(\widehat c_{if}\) is the fitted concentration, while \(\phi_{i\ell}\) is the corresponding feature entry. Before grouping the five training contributions into norms, their scalar form is
\begin{equation}
\equationname{Scalar expansion of the structured-regression training objective}
\label{eq:icsor_training_scalar}
\begin{aligned}
J={}&\sum_{i=1}^{N}\sum_{f=1}^{F}
\left(c_{out,if}-\widehat c_{if}\right)^2\\
&+\lambda_{inv}\sum_{i=1}^{N}\sum_{k=1}^{K}
\left[\sum_{f=1}^{F}A_{kf}(\widehat c_{if}-c_{in,if})\right]^2\\
&+\lambda_{sys}\sum_{i=1}^{N}\sum_{f=1}^{F}
\left[\widehat c_{if}-\sum_{j=1}^{F}\Gamma_{fj}\widehat c_{ij}
-\sum_{\ell=1}^{D}B_{f\ell}\phi_{i\ell}\right]^2\\
&+\lambda_B\sum_{f=1}^{F}\sum_{\ell=1}^{D}B_{f\ell}^2
+\lambda_{\Gamma}\sum_{f=1}^{F}\sum_{j=1}^{F}\Gamma_{fj}^2.
\end{aligned}
\end{equation}
The second contribution first computes an inventory residual by adding component changes with their invariant weights. It then squares that residual. The third contribution compares the fitted component state, after coupling, with the driver generated by the inputs. The last two sums penalize coefficient magnitude rather than prediction residuals.

The dimensions of the matrices show how these sums are assembled. The entry in row \(i\), column \(k\) of \((\widehat C-C_{in})A^\top\) is the inventory residual in the second sum. The entry in row \(i\), column \(f\) of \(\widehat C R^\top\) is \(\widehat c_{if}-\sum_j\Gamma_{fj}\widehat c_{ij}\). The transpose appears because observations are stored as rows, while the single-observation equation \(Rc=r\) uses column vectors.

For \(N\) training observations, \(\Phi\in\mathbb{R}^{N\times D}\) contains feature vectors as rows. The matrices \(C_{in}\) and \(C_{out}\) contain influent and reference effluent states in the same row convention. An auxiliary fitted-state matrix \(\widehat C\in\mathbb{R}_{+}^{N\times F}\) separates the predicted training states from the regression coefficients. The training criterion is
\begin{equation}
\equationname{Matrix form of the structured-regression training objective}
\label{eq:icsor_training_objective}
\begin{aligned}
J(B,\Gamma,\widehat C)={}&
\|C_{out}-\widehat C\|_F^2
+\lambda_{inv}\| (\widehat C-C_{in})A^\top\|_F^2\\
&+\lambda_{sys}\|\widehat C(I_F-\Gamma)^\top-\Phi B^\top\|_F^2\\
&+\lambda_B\|B\|_F^2+\lambda_{\Gamma}\|\Gamma\|_F^2.
\end{aligned}
\end{equation}
The estimation problem minimizes this criterion subject to non-negative fitted training states, zero coupling diagonal, coupling box bounds, and the conditioning safeguard. All penalty weights are non-negative. The adopted ridge weights and coupled-system weight are strictly positive.

The first term measures component prediction error. The second term discourages disagreement with the modeled invariants. The third term encourages agreement between the fitted states and the coupled regression equation. The final two terms limit coefficient magnitude. The non-negative fitted training state balances prediction error, invariant mismatch, and coupled-system mismatch through these weighted terms.

An illustrative one-observation calculation makes the five contributions concrete. Take \(c_{out}=(3,1)^\top\), \(c_{in}=(2,2)^\top\), \(\widehat c=(2.5,1.2)^\top\), and \(A=[1\ \ 1]\). Use the hypothetical coupling matrix in Equation~\eqref{eq:icsor_small_coupling_matrices}. For a single baseline feature, take \(B=(2,1)^\top\), so the driver is \((2,1)^\top\). The five unweighted contributions are
\begin{equation}
\equationname{Illustrative numerical contributions to the training penalties}
\label{eq:icsor_penalty_numbers}
\begin{aligned}
\text{prediction error}&=(3-2.5)^2+(1-1.2)^2=0.29,\\
\text{inventory error}&=(2.5+1.2-2-2)^2=0.09,\\
\text{system error}&=(2.5-0.2(1.2)-2)^2\\
&\quad +(1.2-0.2(2.5)-1)^2=0.1576,\\
\text{driver squared sum}&=2^2+1^2=5,\\
\text{coupling squared sum}&=0.2^2+0.2^2=0.08.
\end{aligned}
\end{equation}
With illustrative weights \(\lambda_{inv}=2\), \(\lambda_{sys}=3\), \(\lambda_B=0.1\), and \(\lambda_{\Gamma}=0.5\), the objective is \(0.29+2(0.09)+3(0.1576)+0.1(5)+0.5(0.08)=1.4828\). These hypothetical weights illustrate the relative cost assigned to each discrepancy. Table~\ref{tab:icsor_fitting_settings} gives the adopted numerical fitting weights.

The auxiliary state is a variable in the fitting problem. The raw prediction for a new input is obtained by solving the learned coupled relation. Residual penalties guide this fit using the general penalty principle of physics-informed training developed by \citet{Raissi2019}. Checked deployment then enforces mass conservation and non-negativity on the returned state.

The objective is expressed in physical component coordinates. Components have different units and numerical scales, so the magnitudes of their residuals and coefficients reflect those conventions. The four penalty weights also depend on the fixed feature and invariant scaling. Consistent rescaling transforms the features, coefficients, invariant operator, reporting map, and penalty weights together. \citet{Pinheiro2025} discuss the influence of feature scaling on learning. For physical interpretation, unit consistency is as necessary as numerical conditioning.

\subsubsection{Deterministic block updates}

The term \(\widehat C\Gamma^\top\) makes the joint objective nonconvex. Holding two variable blocks fixed gives a simpler problem in the remaining block. The fitting procedure updates the driver, coupling matrix, and fitted states in that order within each cycle.

This distinction can be checked with the reduced residual \((1-\gamma z)^2\), where \(\gamma\) is a coupling and \(z\) a fitted component. If \(z\) is fixed, this is a quadratic function of \(\gamma\). If \(\gamma\) is fixed, it is a quadratic function of \(z\). It is not jointly convex. At \((\gamma,z)=(0.2,5)\) and \((0.1,10)\), the residual is zero. At the midpoint \((0.15,7.5)\), it is \((1-1.125)^2=1/64\). A convex function could not exceed the average of the two zero endpoint values at their midpoint. The product between a fitted state and a coupling coefficient accounts for the nonconvex joint objective.

For fixed \(\Gamma\) and \(\widehat C\), the driver update is a ridge-regularized linear least-squares problem. Define \(T=\widehat C R^\top\). For one output row, write its coefficient column as \(\beta\) and its target column from \(T\) as \(t\). The part of the objective that depends on this row is
\begin{equation}
\equationname{Regularized least-squares objective for one driver coefficient row}
\label{eq:icsor_ridge_row_objective}
\lambda_{sys}\sum_{i=1}^{N}
\left(t_i-\sum_{\ell=1}^{D}\phi_{i\ell}\beta_\ell\right)^2
+\lambda_B\sum_{\ell=1}^{D}\beta_\ell^2.
\end{equation}
Differentiate with respect to one coefficient \(\beta_h\). The derivative of the squared residual contributes its residual times the feature \(\phi_{ih}\), while the ridge term contributes \(2\lambda_B\beta_h\). Setting the derivative to zero gives
\begin{equation}
\equationname{Scalar stationarity conditions for a driver coefficient row}
\label{eq:icsor_ridge_scalar_normal}
\lambda_{sys}\sum_{i=1}^{N}\phi_{ih}
\left(\sum_{\ell=1}^{D}\phi_{i\ell}\beta_\ell-t_i\right)
+\lambda_B\beta_h=0.
\end{equation}
Writing this equation for every coefficient gives
\begin{equation}
\equationname{Normal equations for the driver coefficient row}
\label{eq:icsor_ridge_normal_equations}
\left(\lambda_{sys}\Phi^\top\Phi+\lambda_BI_D\right)\beta
=\lambda_{sys}\Phi^\top t.
\end{equation}
The matrix \(\Phi^\top\Phi\) contains sums of products of feature columns. The vector \(\Phi^\top t\) contains sums of feature times target. Adding the positive ridge weight to the diagonal prevents a zero-curvature direction when feature columns are duplicated.

For a hypothetical calculation with one coefficient and no intercept, let \(\Phi=(1,2)^\top\), \(t=(2,3)^\top\), and both weights equal one. The objective becomes
\begin{equation}
\equationname{Illustrative scalar ridge objective and its polynomial expansion}
\label{eq:icsor_ridge_numeric_example}
(2-\beta)^2+(3-2\beta)^2+\beta^2
=13-16\beta+6\beta^2.
\end{equation}
Its derivative is \(-16+12\beta\), so the minimizer is \(\beta=4/3\). The matrix calculation gives \(\Phi^\top\Phi=1^2+2^2=5\), \(\Phi^\top t=1(2)+2(3)=8\), and \((5+1)\beta=8\). Both calculations agree. Without the penalty, the coefficient would be \(8/5\). The reduction illustrates ridge shrinkage for one coefficient. The adopted driver still includes its baseline and all stated features.

Assembling the normal equations for all output rows gives the exact solution
\begin{equation}
\equationname{Closed-form ridge update for the driver coefficient matrix}
\label{eq:icsor_ridge_update}
B^\top=
\left(\lambda_{sys}\Phi^\top\Phi+\lambda_B I_D\right)^{-1}
\lambda_{sys}\Phi^\top\widehat C R^\top.
\end{equation}
The positive driver ridge weight makes the matrix in this expression positive definite, even when ordered quadratic columns are duplicated. In computation, the corresponding linear system is solved. An explicit matrix inverse is unnecessary. The square quadratic slices are then symmetrized without changing their predictions.

For fixed \(B\) and \(\widehat C\), define \(E=\widehat C-\Phi B^\top\).

The coupling calculation can be read one output row at a time. For row \(f\), the residual at observation \(i\) is \(E_{if}-\sum_{j\neq f}\widehat c_{ij}\Gamma_{fj}\). Its minimization uses the other fitted component columns as predictors of \(E_{if}\), with a ridge penalty and box bounds. In a hypothetical row with one permitted coupling \(\gamma\), two predictor values 1 and 2, and residual targets 0.4 and 0.8, equal unit weights give
\begin{equation}
\equationname{Illustrative scalar coupling objective and its expansion}
\label{eq:icsor_coupling_numeric_update}
(0.4-\gamma)^2+(0.8-2\gamma)^2+\gamma^2
=0.8-4\gamma+6\gamma^2.
\end{equation}
The derivative \(-4+12\gamma\) vanishes at \(\gamma=1/3\). This value lies below the adopted absolute entry bound 0.368. If the illustrative box bound were instead 0.25, the convex scalar problem would attain its minimum at 0.25. The conditioning test is applied to the complete coupling matrix after these row-wise calculations.

Collecting these scalar residuals across all output rows gives the coupling proposal
\begin{equation}
\equationname{Regularized least-squares update for component coupling}
\label{eq:icsor_coupling_update}
\min_{\Gamma}\quad
\lambda_{sys}\|E-\widehat C\Gamma^\top\|_F^2+
\lambda_{\Gamma}\|\Gamma\|_F^2
\end{equation}
subject to its zero diagonal and box bounds. This is a convex quadratic program. Its candidate is subsequently screened using Equation~\eqref{eq:icsor_coupling_safeguards}. A candidate that fails the condition check is shrunk toward zero or replaced by an admissible fallback. The screened proposal is evaluated against the previous admissible coupling matrix using the complete objective.

For fixed \(B\) and \(\Gamma\), the fitted-state update separates over the training rows. For sample \(i\), write \(\widehat c_i\) as a column vector.

Expand the three residual sums for one observation before collecting their coefficients. The data term contributes \(\widehat c^\top\widehat c-2c_{out}^\top\widehat c\). The invariant term contributes \(\lambda_{inv}\widehat c^\top A^\top A\widehat c-2\lambda_{inv}c_{in}^\top A^\top A\widehat c\). The system term contributes \(\lambda_{sys}\widehat c^\top R^\top R\widehat c-2\lambda_{sys}(B\phi)^\top R\widehat c\). Terms independent of \(\widehat c\) do not affect its minimizer. Collecting its quadratic terms gives \(H_C\), and collecting its linear terms gives \(h_i\).

Write the expanded quadratic as \(\widehat c_i^\top H_C\widehat c_i-2h_i^\top\widehat c_i\). The entry \((H_C)_{fg}\) multiplies the ordered product \(\widehat c_{if}\widehat c_{ig}\), and \((h_i)_f\) supplies its linear term. For distinct components, the two mirrored entries contribute together, following the quadratic symmetry convention. Let \(\delta_{fg}\) equal one when \(f=g\) and zero otherwise. Reading the coefficients entry by entry gives
\begin{equation}
\equationname{Scalar entries of the fitted-state Hessian and forcing vector}
\label{eq:icsor_state_scalar_coefficients}
\begin{aligned}
(H_C)_{fg}&=\delta_{fg}+\lambda_{inv}\sum_{k=1}^{K}A_{kf}A_{kg}
+\lambda_{sys}\sum_{j=1}^{F}R_{jf}R_{jg},\\
(h_i)_f&=c_{out,if}
+\lambda_{inv}\sum_{k=1}^{K}A_{kf}\sum_{g=1}^{F}A_{kg}c_{in,ig}\\
&\quad+\lambda_{sys}\sum_{j=1}^{F}R_{jf}\sum_{\ell=1}^{D}B_{j\ell}\phi_{i\ell}.
\end{aligned}
\end{equation}
For example, the invariant contribution to \((H_C)_{12}\) adds the product of the first and second component weights over the invariant rows. The system contribution adds the product of columns 1 and 2 of \(R\). These products explain why correcting one fitted component can change the preferred value of another.

Completing the square explains what the non-negative solve is adjusting. With \(c_0=H_C^{-1}h_i\),
\begin{equation}
\equationname{Completed-square form of the fitted-state objective}
\label{eq:icsor_state_completed_square}
\widehat c^\top H_C\widehat c-2h_i^\top\widehat c
=(\widehat c-c_0)^\top H_C(\widehat c-c_0)
-h_i^\top H_C^{-1}h_i.
\end{equation}
The unrestricted minimizer is \(c_0\). When it contains a negative entry, the constrained problem seeks the closest non-negative state in the quadratic geometry defined by \(H_C\). Off-diagonal entries tilt that geometry, so ordinary coordinate clipping need not give the closest state.

For example, take hypothetical \(H_C=\begin{bmatrix}3&1\\1&3\end{bmatrix}\) and \(h_i=(-1,2)^\top\). These coefficients can arise from unit data, invariant, and system weights with \(A=[1\ \ 1]\) and \(R=I\). A signed driver can make the first entry of \(h_i\) negative. The unrestricted solution is \(c_0=(-5/8,7/8)^\top\). Clipping gives \((0,7/8)^\top\). Instead, set the constrained first component to zero and minimize the remaining expression \(3c_2^2-4c_2\). Its derivative \(6c_2-4\) gives \(c_2=2/3\).

At \((0,2/3)^\top\), the derivative in the first coordinate is \(2(c_2+1)=10/3\), which is positive. Increasing that zero component would increase the objective, so the boundary value is appropriate. The objective is \(-4/3\) at this state and \(-77/64\) at the clipped state. The constrained solution has the smaller objective. The negative objective values are possible because the constants removed during expansion are absent. The complete residual-based training criterion remains non-negative.

The general coefficient arrays derived above are assembled in the compact form
\begin{equation}
\equationname{Hessian and forcing vector for the fitted-state update}
\label{eq:icsor_state_update}
\begin{aligned}
H_C&=I_F+\lambda_{inv}A^\top A+\lambda_{sys}R^\top R,\\
h_i&=c_{out,i}+\lambda_{inv}A^\top A c_{in,i}
+\lambda_{sys}R^\top B\phi_i.
\end{aligned}
\end{equation}
Each row minimizes \(\widehat c_i^\top H_C\widehat c_i-2h_i^\top\widehat c_i\) over \(\widehat c_i\geq\mathbf{0}\). The identity term makes \(H_C\) positive definite, so the row minimizer is unique. Clipping the unconstrained solution componentwise generally does not solve this problem because the off-diagonal terms in \(H_C\) connect component adjustments.

The quadratic-program method of \citet{Stellato2018} provides an algorithmic foundation for the constrained coupling and fitted-state blocks. These blocks use their quadratic structure directly. A generic stochastic first-order method could also optimize a related criterion, but it would require separate mechanisms to control non-negativity, coupling bounds, and conditioning. The block formulation keeps each requirement attached to an explicit subproblem.

\subsubsection{Descent, stopping, and numerical safeguards}

Exact minimization over a block cannot increase the objective when its admissible set contains the previous block value. Since \(J\geq0\), an objective sequence that is monotone non-increasing converges in value. The convergence analysis of \citet{Razaviyayn2013} relates blockwise minimization to stationary solutions under its stated conditions.

The conditioning screen requires particular care. Shrinking a coupling proposal can increase the training criterion relative to the previous admissible coupling matrix. The procedure therefore evaluates the objective after screening and retains the previous admissible matrix when a screened candidate cannot provide an acceptable non-increasing update. Solver tolerances also require an explicit allowance for numerical error. Monotone descent is assessed for the accepted objective-decreasing updates after screening.

Each restart begins with an admissible coupling matrix and non-negative fitted states. A ridge solve supplies the initial driver. Full cycles continue until an iteration limit is reached or the slope of the trailing objective history is sufficiently small. A running minimum records the best admissible state reached. When several restarts are used, the admissible state with the smallest recorded objective is retained. A nearly flat objective triggers the prescribed stopping rule.

Figure~\ref{fig:ch5_training_cycle} separates the three fitting updates from the checks that make an update acceptable. The scalar calculations in Equations~\eqref{eq:icsor_ridge_numeric_example} and \eqref{eq:icsor_coupling_numeric_update} illustrate the first and third computational boxes in the left column. The non-negative-state example explains the return path in the right column.

\begin{figure}[htbp]
\centering
\input{figures/ch5_concept_training_cycle.tex}
\caption{Conceptual fitting cycle with explicit symmetry, conditioning, descent, and non-negativity checks. The prescribed restart count can be one.}
\label{fig:ch5_training_cycle}
\end{figure}

The cycle returns to the driver because each block changes the conditions faced by the others. The conditioning screen acts on the coupling proposal before the fitted-state update. The full objective is then recorded after the cycle. The restart decision concerns the prescribed search procedure, while the final selection compares admissible objective values. These steps define the accepted sequence of updates for the nonconvex joint fit.

\begin{table}[htbp]
\centering
\caption{Specified fitting settings for the structured regression.}
\label{tab:icsor_fitting_settings}
\begin{tabular}{ll}
\toprule
Setting & Value\\
\midrule
Baseline feature & Included\\
Invariant penalty \(\lambda_{inv}\) & \(8.68\times10^{-3}\)\\
Coupled-system penalty \(\lambda_{sys}\) & \(1.27\times10^{-3}\)\\
Driver ridge penalty \(\lambda_B\) & \(6.09\times10^{-5}\)\\
Coupling ridge penalty \(\lambda_{\Gamma}\) & \(6.55\times10^{-4}\)\\
Off-diagonal bound \(\gamma_{abs}\) & 0.368\\
Maximum full fitting cycles & 60\\
Restarts & 1\\
Trailing objective window & 12 cycles\\
Objective slope tolerance & \(2.85\times10^{-7}\)\\
\bottomrule
\end{tabular}
\end{table}

Table~\ref{tab:icsor_fitting_settings} specifies numerical fitting choices. The condition threshold \(\kappa_{max}\) remains an additional prescribed numerical safeguard and is checked independently of the coefficient bound. Hyperparameter assessment uses validation observations within the training pool. Held-out test observations do not determine penalty weights or stopping choices. Fitting settings remain fixed within a sample-size comparison so that changes in the learning curve reflect changing sample availability under a common model specification.

\subsection{Checked deployment}

\subsubsection{The raw coupled state}

For a new pair \((u_*,c_{in,*})\), the fitted parameters give
\begin{equation}
\equationname{Raw deployed state from the fitted coupled system}
\label{eq:icsor_raw_state}
\widehat R=I_F-\widehat\Gamma,
\qquad d_* =\widehat B\phi(u_*,c_{in,*}),
\qquad \widehat R c_{raw}=d_*.
\end{equation}
The notation \(c_{raw}=\widehat R^{-1}d_*\) describes the mapping mathematically. A linear-system solve evaluates it numerically. The accepted conditioning bound limits the sensitivity of that solve. The raw state is returned only when it satisfies both the invariant and non-negativity checks.

For a vector, the Euclidean norm is the square root of the sum of squared entries. Thus an illustrative invariant-residual vector \((3,4)^\top\) has norm \(\sqrt{3^2+4^2}=5\). The absolute norm adds the absolute entries instead. For an illustrative concentration vector \((-0.3,0.1,-0.4)^\top\), retaining only negative entries gives \((-0.3,0,-0.4)^\top\), whose absolute norm is 0.7. The first calculation measures the size of a mass conservation discrepancy. The second measures total negative concentration in the stated numerical coordinates.

For any candidate \(c\), define the diagnostic residuals
\begin{equation}
\equationname{Mass conservation and non-negativity diagnostics of a deployed state}
\label{eq:icsor_deployment_diagnostics}
v_{mc}(c)=\|A(c-c_{in,*})\|_2,
\qquad
v_{nn}(c)=\|\min(c,\mathbf{0})\|_1.
\end{equation}
The minimum is taken componentwise. The adopted diagnostic tolerance is \(10^{-10}\). Algebraic equality and non-negativity define the mathematical requirement. Residuals below the numerical tolerance define its computational assessment. Invariant residual magnitudes use the stated scaling of \(A\).

\subsubsection{Affine projection for mass conservation}

The nearest state satisfying mass conservation can first be derived for the hypothetical two-component inventory. Write the raw prediction as \((p_1,p_2)^\top\) and the required total as \(T\). The two unknown projected concentrations minimize \((c_1-p_1)^2+(c_2-p_2)^2\) while satisfying \(c_1+c_2=T\). Introduce a multiplier \(\eta\) for that equality and write
\begin{equation}
\equationname{Lagrangian for an illustrative two-component affine projection}
\label{eq:icsor_scalar_projection_lagrangian}
\mathcal{L}=\frac{1}{2}(c_1-p_1)^2+
\frac{1}{2}(c_2-p_2)^2+\eta(c_1+c_2-T).
\end{equation}
The factor one-half simplifies differentiation without changing the minimizer. Differentiating with respect to the two concentrations and the multiplier gives
\begin{equation}
\equationname{Stationarity and equality conditions for the illustrative affine projection}
\label{eq:icsor_scalar_projection_conditions}
c_1-p_1+\eta=0,
\qquad c_2-p_2+\eta=0,
\qquad c_1+c_2=T.
\end{equation}
The first two equations give \(c_1=p_1-\eta\) and \(c_2=p_2-\eta\). Substituting them into the inventory equality gives \(p_1+p_2-2\eta=T\). Thus
\begin{equation}
\equationname{Closed-form solution of the illustrative affine projection}
\label{eq:icsor_scalar_projection_solution}
\begin{aligned}
\eta&=\frac{p_1+p_2-T}{2},\\
c_1&=p_1-\frac{p_1+p_2-T}{2},\\
c_2&=p_2-\frac{p_1+p_2-T}{2}.
\end{aligned}
\end{equation}
The unweighted distance shares the inventory discrepancy equally between these two coordinates. With \(p_1=12\), \(p_2=1\), and \(T=10\), the discrepancy is 3 and the multiplier is 1.5. The projected state is \((10.5,-0.5)^\top\). The equality is satisfied even though the second component is negative. This example identifies the inequality that is absent from an equality-only projection.

When the raw state fails its checks, the first projection solves
\begin{equation}
\equationname{Euclidean affine projection onto the mass conservation equality}
\label{eq:icsor_affine_problem}
c_{aff}=\mathop{\arg\min}_{c\in\mathbb{R}^{F}}\|c-c_{raw}\|_2^2
\quad\text{subject to}\quad Ac=Ac_{in,*}.
\end{equation}
Introduce a Lagrange multiplier \(\eta\in\mathbb{R}^{K}\). Stationarity gives \(c-c_{raw}+A^\top\eta=0\). Substitution into the equality gives \(AA^\top\eta=A(c_{raw}-c_{in,*})\). The full row rank of \(A\) therefore gives the unique solution
\begin{equation}
\equationname{Closed-form affine projection of the raw component state}
\label{eq:icsor_affine_solution}
c_{aff}=c_{raw}-A^\top(AA^\top)^{-1}A(c_{raw}-c_{in,*}).
\end{equation}
The fixed matrix \(A^\top(AA^\top)^{-1}A\) can be prepared once. The affine projection is then a matrix-vector operation for each new prediction.

Each matrix step follows the same scalar derivation. The multiplier column contains one multiplier per independent inventory. The vector \(A^\top\eta\) converts these inventory projections into component adjustments. Solving \(AA^\top\eta=A(c_{raw}-c_{in,*})\) determines the adjustments needed to remove the inventory residual. For the two-component illustration,
\begin{equation}
\equationname{Illustrative matrix of the mass conservation projection}
\label{eq:icsor_small_affine_matrix}
AA^\top=2,
\qquad
A^\top(AA^\top)^{-1}A
=\frac{1}{2}\begin{bmatrix}1&1\\1&1\end{bmatrix}.
\end{equation}
Taking the illustrative influent as \((6,4)^\top\), the raw-minus-influent difference is \((6,-3)^\top\). Multiplying it by this projection matrix gives \((1.5,1.5)^\top\), which is subtracted from the raw state. This recovers the scalar calculation without guessing how much either component should change.

This projection is an orthogonal projection onto the affine set defined by mass conservation in the chosen numerical component coordinates. It removes the part of the prediction error normal to that set and retains the part tangent to it. The projection distance uses the stated component unit convention. \citet{SturmSilva2024} develop related projections for mass conservation and explain why species-dependent weights matter when concentration scales differ.

The affine state is checked for both mass conservation residuals and non-negativity before deployment. The equality follows analytically from Equation~\eqref{eq:icsor_affine_solution}, but checking its residual catches numerical or operator inconsistency. An affine state that satisfies mass conservation can still contain a negative concentration. That possibility requires a projection with inequality constraints.

\subsubsection{Joint mass conservation and non-negativity}

The absolute-distance projection can also be written before introducing compact vector inequalities. For the state that satisfies mass conservation but violates non-negativity \((10.5,-0.5)^\top\), require \(c_1+c_2=10\) and \(c_1,c_2\geq0\). Introduce non-negative adjustment bounds \(\rho_1,\rho_2\). The constraints
\begin{equation}
\equationname{Absolute-adjustment bounds for an illustrative two-component projection}
\label{eq:icsor_scalar_absolute_constraints}
\begin{aligned}
\rho_1&\geq c_1-10.5,&\rho_1&\geq10.5-c_1,\\
\rho_2&\geq c_2+0.5,&\rho_2&\geq-0.5-c_2
\end{aligned}
\end{equation}
make \(\rho_1\) at least the absolute first adjustment and \(\rho_2\) at least the absolute second adjustment. Minimizing \(w_1\rho_1+w_2\rho_2\) with positive weights makes each bound equal to the corresponding absolute adjustment. No absolute-value function remains in the optimization constraints.

Set \(c_2=t\) to impose the inventory equality, so \(c_1=10-t\) and \(0\leq t\leq10\). Both absolute adjustments then equal \(t+0.5\). The objective is \((w_1+w_2)(t+0.5)\), which increases with \(t\). Its minimum is at \(t=0\), giving \((10,0)^\top\). For equal unit weights, its cost is 1. At \(t=1\), the state would be \((9,1)^\top\), with cost 3. This direct comparison explains the optimal projection in the reduced problem.

When the affine state fails the non-negativity check, the final projection solves
\begin{equation}
\equationname{Weighted absolute-distance projection for mass conservation and non-negativity}
\label{eq:icsor_weighted_correction}
\begin{aligned}
\widehat c =\mathop{\arg\min}_{c,\rho\in\mathbb{R}^{F}}\quad
& w_c^\top\rho\\
\text{subject to}\quad
& Ac=Ac_{in,*},\\
& c\geq\mathbf{0},\\
& -\rho\leq c-c_{aff}\leq\rho,\\
& \rho\geq\mathbf{0}.
\end{aligned}
\end{equation}
For strictly positive fixed weights \(w_c\), minimization makes \(\rho_f=|c_f-c_{aff,f}|\) at an optimum. The objective is therefore the weighted absolute adjustment \(\sum_f w_{c,f}|c_f-c_{aff,f}|\). The program minimizes adjustment from the affine state while retaining the fitted regression coefficients.

The program is linear because the absolute values are represented by auxiliary variables and linear inequalities. The revised-simplex method described by \citet{HuangfuHall2018} provides an algorithmic foundation for solving this type of projection. The mathematical feasible set is nonempty whenever the supplied influent state is non-negative under the adopted boundary relation. The positive weights make the adjustment objective bounded below and discourage unconstrained movement of any component. A minimizer exists, although several minimizing states can exist when the absolute-distance objective has a flat face.

Large weights make changes to a component more costly. Small weights permit larger changes when those changes help satisfy the specified constraints. Weights are fixed before deployment using the intended component scaling and scientific priorities. A sensitivity analysis of the weights can examine how much the returned state depends on that modeling choice. Equal numerical weights also embody a choice because component units and typical magnitudes differ.

Three hypothetical components show how weights can change an admissible allocation. Let the inventory row be \([1\ \ 1\ \ 1]\) and the affine reference satisfying mass conservation be \((-1,4,7)^\top\), whose total is 10. Restoring its first component to zero requires removal of one unit from another component. Two feasible candidates are \((0,3,7)^\top\) and \((0,4,6)^\top\). For weights \((1,10,1)^\top\), their costs are respectively \(1(1)+10(1)=11\) and \(1(1)+1(1)=2\). The second component is more costly to change, so removing material from the third is preferable. With weights \((1,1,10)^\top\), the two costs become 2 and 11. The preferred allocation reverses. Both allocations preserve the same inventory, while the weights determine the preferred adjustment.

\citet{KircherVotsmeier2025} show why a projection for mass conservation needs simultaneous non-negativity enforcement in reaction-kinetics surrogates. Their scaled-space treatment protects low-concentration species during projection. Equation~\eqref{eq:icsor_weighted_correction} addresses the same joint requirement through a weighted absolute-distance problem for activated sludge components. It imposes the adopted linear invariants and non-negativity on the returned state.

\subsubsection{Illustrative projection and the failure of clipping}

Consider an illustrative two-component system with \(A=[1\ \ 1]\), a total influent inventory of 10, and a raw prediction \(c_{raw}=(12,1)^\top\). The raw total is 13. The affine formula subtracts 1.5 from each component and gives \(c_{aff}=(10.5,-0.5)^\top\). The total is now 10, but the second component is negative.

Clipping that negative value to zero produces \((10.5,0)^\top\), which violates the total inventory. With equal projection weights, Equation~\eqref{eq:icsor_weighted_correction} instead returns \((10,0)^\top\). Both constraints hold. These hypothetical values illustrate the projection geometry.

The deployment order uses this geometry directly. A feasible raw state needs no adjustment. An infeasible raw state is first projected onto the mass conservation set. A feasible affine state can then be returned. The linear program is required only when the affine state still fails non-negativity. Every returned candidate is checked under the declared numerical tolerances. An unsuccessful numerical solve is recorded as a prediction failure.

Figure~\ref{fig:ch5_checked_deployment} makes the branch conditions explicit. Both physical checks appear at the raw and affine stages. The last projection branch also requires successful optimization and an independent admissibility check.

\begin{figure}[htbp]
\centering
\input{figures/ch5_concept_checked_deployment.tex}
\caption{Checked deployment of one predicted component state. Accepted branches merge before the common reporting map is applied. A failed final audit produces a reported prediction failure.}
\label{fig:ch5_checked_deployment}
\end{figure}

The hypothetical state \((12,1)^\top\) follows the two No branches because its total is incorrect and its affine projection contains a negative component. The weighted projection supplies \((10,0)^\top\), which reaches the return branch only after the stated checks pass. A state that satisfies both conditions at an earlier branch reaches the same reporting step. The different paths therefore change the projection effort while retaining the same acceptance requirements.

The component state returned by this rule is the deployed output. Its reported quantities are
\begin{equation}
\equationname{Water-quality composites of the final checked state}
\label{eq:icsor_deployed_composites}
\widehat y_{\mathrm{comp}}=I_{\mathrm{comp}}\widehat c.
\end{equation}
The deployment constraints enforce admissibility, and the stage-wise assessment measures how often each projection is needed.

\subsection{Interpreting the complete prediction map}

\subsubsection{Driver coefficients and back projection}

The transformation from driver coefficients to output coefficients can be computed with small matrices. Use the hypothetical coupling matrix in Equation~\eqref{eq:icsor_small_coupling_matrices}, a reduced feature vector \((1,a)^\top\), and driver coefficients
\begin{equation}
\equationname{Illustrative driver coefficients and operating-dependent driver}
\label{eq:icsor_small_driver_coefficient_matrix}
B=\begin{bmatrix}4&1\\2&0\end{bmatrix},
\qquad
r=\begin{bmatrix}4+a\\2\end{bmatrix}.
\end{equation}
The first coefficient column describes the baseline, and the second describes dependence on the hypothetical scalar input \(a\). Multiplication through the coupling inverse gives
\begin{equation}
\equationname{Illustrative transformation of driver coefficients through coupling}
\label{eq:icsor_small_back_projection}
\begin{aligned}
R^{-1}B
&=\frac{1}{0.96}
\begin{bmatrix}1&0.2\\0.2&1\end{bmatrix}
\begin{bmatrix}4&1\\2&0\end{bmatrix}\\
&=\frac{1}{0.96}\begin{bmatrix}4.4&1\\2.8&0.2\end{bmatrix}
=\begin{bmatrix}55/12&25/24\\35/12&5/24\end{bmatrix}.
\end{aligned}
\end{equation}
The second component now depends on \(a\), even though its driver coefficient for \(a\) was zero. This dependence is introduced by the simultaneous coupling solve.

For a hypothetical reported quantity \(y=c_1+2c_2\), its reporting row is \([1\ \ 2]\). Its transformed coefficient row is
\begin{equation}
\equationname{Illustrative composite coefficients after coupling}
\label{eq:icsor_small_composite_projection}
\begin{bmatrix}1&2\end{bmatrix}R^{-1}B
=\begin{bmatrix}125/12&35/24\end{bmatrix}.
\end{equation}
At \(a=2\), the driver is \((6,2)^\top\), the raw component state is \((20/3,10/3)^\top\), and the reported quantity is \(40/3\). The transformed row gives the same value through \(125/12+(35/24)2=40/3\). The hypothetical reporting row provides a reduced arithmetic example.

The rows of \(\widehat B\) refer to component drivers. Coupling and the reporting map convert these rows into COD, TN, TP, and TSS coefficients. The raw component coefficient matrix is
\begin{equation}
\equationname{Raw component-response coefficients after coupling}
\label{eq:icsor_raw_component_coefficients}
B_{\mathrm{raw}}=\widehat R^{-1}\widehat B,
\qquad c_{raw}=B_{\mathrm{raw}}\phi.
\end{equation}
The corresponding map in reported-composite space is
\begin{equation}
\equationname{Raw composite-response coefficients after coupling and reporting}
\label{eq:icsor_raw_composite_coefficients}
B_{\mathrm{raw,comp}}=I_{\mathrm{comp}}\widehat R^{-1}\widehat B,
\qquad y_{\mathrm{raw,comp}}=B_{\mathrm{raw,comp}}\phi.
\end{equation}
This transformation carries the driver coefficients through the learned coupling and then through the reporting operator. Inspection of a COD quadratic block uses its row in \(B_{\mathrm{raw,comp}}\), rather than an arbitrary row of \(\widehat B\).

Figure~\ref{fig:ch5_prediction_chain} shows the complete raw prediction chain using the named feature groups. It keeps the component driver, simultaneous coupling solve, and reporting operation visible as separate mathematical steps.

\begin{figure}[htbp]
\centering
\input{figures/ch5_concept_prediction_chain.tex}
\caption{Named feature blocks and the transformation from driver to raw component and reported-composite predictions. Composing the three operations gives the raw reported coefficient map \(I_{\mathrm{comp}}R^{-1}B\).}
\label{fig:ch5_prediction_chain}
\end{figure}

In the two-feature example, the second driver has no direct coefficient for the scalar input, but the coupling path gives its raw component a nonzero input coefficient. The reporting path combines both component coefficients. Reading the paths in this order explains the back projection in Equations~\eqref{eq:icsor_small_back_projection} and \eqref{eq:icsor_small_composite_projection}. Checked deployment acts on the raw component state before the final reported quantities are calculated.

The named blocks can be displayed as coefficient maps. Operating terms are grouped separately from influent terms. Symmetric quadratic maps display both ordered cells while reporting unique cross terms consistently. Operating-load maps show which operating variables interact with which influent components. The coupling matrix is displayed separately because it has a different interpretation from the input-driver terms.

A coefficient-display threshold is used only to simplify reporting after fitting. For each reported-composite raw-regression map, the threshold is \(10^{-4}\) times its largest absolute coefficient. For the coupling matrix, the threshold is \(10^{-4}\) times its largest absolute off-diagonal entry. The fitted mapping is retained in full, and the threshold determines which entries appear in the coefficient display. Symmetric quadratic summaries count unique upper-triangular terms. A small coefficient can still have a large contribution when multiplied by a large feature. Interpretation therefore examines coefficient magnitude, input scale, and the contribution of the corresponding term over the declared domain.

\subsubsection{Analytical jacobians of the raw regression}

Differentiation of a scalar polynomial shows why a coefficient and a marginal response differ. For the illustrative driver in Equation~\eqref{eq:icsor_scalar_driver_example},
\begin{equation}
\equationname{Gradients of the illustrative scalar regression driver}
\label{eq:icsor_scalar_driver_derivatives}
\begin{aligned}
\frac{\partial r_1}{\partial u_1}&=2+0.2u_1+0.2u_2,\\
\frac{\partial r_1}{\partial u_2}&=-1+0.2u_1+0.25c_1,\\
\frac{\partial r_1}{\partial c_1}&=0.5+0.25u_2,
&\frac{\partial r_1}{\partial c_2}&=-0.1c_2.
\end{aligned}
\end{equation}
At \(u=(2,3)^\top\) and \(c=(4,5)^\top\), these derivatives are 3, 0.4, 1.25, and \(-0.5\). In particular, the first derivative is 3 although the linear coefficient of \(u_1\) is 2. The operating quadratic term and interaction with \(u_2\) supply the remaining contribution.

The derivative of the second illustrative driver \(r_2=-u_1+u_2+0.2c_2\) with respect to \(u_1\) is \(-1\). Differentiate the coupled equations with the coupling held fixed. Their derivative equations are \(\partial c_1/\partial u_1-0.2\partial c_2/\partial u_1=3\) and \(-0.2\partial c_1/\partial u_1+\partial c_2/\partial u_1=-1\). Thus the same inverse used for prediction gives
\begin{equation}
\equationname{Illustrative propagation of an operating derivative through coupling}
\label{eq:icsor_small_chain_derivative}
\begin{bmatrix}\partial c_1/\partial u_1\\\partial c_2/\partial u_1\end{bmatrix}
=\frac{1}{0.96}\begin{bmatrix}1&0.2\\0.2&1\end{bmatrix}
\begin{bmatrix}3\\-1\end{bmatrix}
=\begin{bmatrix}35/12\\-5/12\end{bmatrix}.
\end{equation}
For the hypothetical reporting row \([1\ \ 2]\), the reported derivative is \(35/12+2(-5/12)=25/12\). Differentiating the driver, solving through the coupling, and applying the reporting weights are three separate operations.

Putting several derivatives in one matrix gives a Jacobian. Its rows identify the output being differentiated, and its columns identify the input being varied. At the illustrative inputs, the operating-driver Jacobian and its coupled transformation are
\begin{equation}
\equationname{Illustrative operating Jacobians before and after coupling}
\label{eq:icsor_small_operating_jacobian}
J_{r,u}=\begin{bmatrix}3&0.4\\-1&1\end{bmatrix},
\qquad
R^{-1}J_{r,u}=\begin{bmatrix}35/12&5/8\\-5/12&9/8\end{bmatrix}.
\end{equation}
The first column is the derivative calculation above. The second follows by solving the same coupled system for the derivative column \((0.4,1)^\top\). A Jacobian therefore collects several scalar sensitivity calculations rather than introducing a different kind of response.

A coefficient describes a term in the polynomial. A Jacobian describes the change in the whole mapping at an operating point. For component \(f\), let \(H_{uc}^{(f)}\in\mathbb{R}^{M_{op}\times F}\) reshape its interaction coefficients so that the cross term is \(u^\top H_{uc}^{(f)}c_{in}\). The driver can then be expressed as
\begin{equation}
\equationname{Component driver expressed through symmetric quadratic matrices}
\label{eq:icsor_driver_component}
\begin{aligned}
r_f={}&b_f+(W_u)_{f,:}u+(W_{in})_{f,:}c_{in}
+u^\top Q_{uu}^{(f)}u\\
&+c_{in}^\top Q_{cc}^{(f)}c_{in}
+u^\top H_{uc}^{(f)}c_{in}.
\end{aligned}
\end{equation}
With symmetric square slices, its gradients are
\begin{equation}
\equationname{Operating and influent gradients of a component driver}
\label{eq:icsor_driver_gradients}
\begin{aligned}
\nabla_u r_f&=(W_u)_{f,:}^\top+2Q_{uu}^{(f)}u
+H_{uc}^{(f)}c_{in},\\
\nabla_{c_{in}}r_f&=(W_{in})_{f,:}^\top+2Q_{cc}^{(f)}c_{in}
+(H_{uc}^{(f)})^\top u.
\end{aligned}
\end{equation}
The quadratic gradients in Equation~\eqref{eq:icsor_driver_gradients} follow by differentiating each scalar product. For \(x^\top Qx=\sum_i\sum_jq_{ij}x_ix_j\), differentiation with respect to \(x_h\) collects the terms with \(i=h\) and those with \(j=h\). Their sum is \(\sum_j(q_{hj}+q_{jh})x_j\). It becomes the \(h\)th entry of \(2Qx\) when \(Q\) is symmetric. For \(u^\top Hc=\sum_i\sum_jh_{ij}u_ic_j\), differentiation with respect to \(u_i\) gives \(\sum_jh_{ij}c_j\), while differentiation with respect to \(c_j\) gives \(\sum_i h_{ij}u_i\). These are the entries of \(Hc\) and \(H^\top u\).

Stacking the transposed gradients gives \(J_{r,u}\in\mathbb{R}^{F\times M_{op}}\) and \(J_{r,c}\in\mathbb{R}^{F\times F}\). The raw-state and raw-composite Jacobians follow through the same back projection as the coefficients,
\begin{equation}
\equationname{Raw component and composite Jacobians after coupling}
\label{eq:icsor_raw_jacobians}
\begin{aligned}
J_{raw,u}&=\widehat R^{-1}J_{r,u},
&J_{raw,c}&=\widehat R^{-1}J_{r,c},\\
J_{raw,comp,u}&=I_{\mathrm{comp}}J_{raw,u},
&J_{raw,comp,c}&=I_{\mathrm{comp}}J_{raw,c}.
\end{aligned}
\end{equation}
These derivatives describe the complete fitted raw regression at the specified inputs, including curvature and interactions. The projection-stage derivatives are calculated below.

An illustrative aeration sensitivity can be inspected at low and high ammonium loading while holding the remaining inputs fixed. A difference in the two derivatives identifies a loading-dependent association within the fitted model. Oxygen-transfer limitations and heterotrophic competition provide process context for assessing such a pattern, as discussed by \citet{StenstromSong1991}. Comparison with that process context provides a plausibility check on the fitted response pattern.

\subsubsection{Sensitivity after the affine projection}

Return to the two-feature example in Equation~\eqref{eq:icsor_small_back_projection}. Its raw derivative with respect to \(a\) is \((25/24,5/24)^\top\), whose total is \(5/4\). This raw response changes the total component inventory. At fixed influent inventory, an affine projection must remove that total change. For the two-component invariant,
\begin{equation}
\equationname{Illustrative operating derivative after affine projection}
\label{eq:icsor_small_projection_derivative}
\begin{aligned}
P_A&=\frac{1}{2}\begin{bmatrix}1&-1\\-1&1\end{bmatrix},\\
\frac{\partial c_{aff}}{\partial a}
&=P_A\begin{bmatrix}25/24\\5/24\end{bmatrix}
=\begin{bmatrix}5/12\\-5/12\end{bmatrix}.
\end{aligned}
\end{equation}
The two entries now sum to zero. For the reporting row \([1\ \ 2]\), the derivative is \(5/12+2(-5/12)=-5/12\). The corresponding raw derivative was \(35/24\). The projection changes both its magnitude and sign because it enforces a different material-accounting condition on the complete returned state.

The influent reference also changes when an influent component is perturbed. In this reduced example, the driver contains no influent term. Increasing either influent component by \(\delta\) increases the required total by \(\delta\). The unweighted affine projection then adds \(\delta/2\) to each predicted component. Its influent derivative is the matrix \(\frac12\begin{bmatrix}1&1\\1&1\end{bmatrix}\), even though the raw influent derivative is zero. This contribution comes from the right-hand side of the mass conservation equality.

Define the complementary projection matrices
\begin{equation}
\equationname{Complementary mass conservation projection matrices}
\label{eq:icsor_projection_matrices}
Q_A=A^\top(AA^\top)^{-1}A,
\qquad P_A=I_F-Q_A.
\end{equation}
Then \(c_{aff}=P_Ac_{raw}+Q_Ac_{in}\). With the fitted coefficients and invariant operator held fixed, its Jacobians are
\begin{equation}
\equationname{Operating and influent Jacobians of the affine-projected state}
\label{eq:icsor_affine_jacobians}
J_{aff,u}=P_AJ_{raw,u},
\qquad J_{aff,c}=P_AJ_{raw,c}+Q_A.
\end{equation}
The second expression includes the change in the mass conservation reference when the influent changes. The term \(Q_A\) accounts for the changing influent reference.

Multiplication by \(A\) gives \(AJ_{aff,u}=0\) and \(AJ_{aff,c}=A\). An operating perturbation changes the affine prediction only along directions that preserve its fixed influent inventories. An influent perturbation changes those inventories consistently with the influent reference. These relations provide analytical checks on the sensitivity calculation.

When the affine branch is used over a neighborhood of inputs, its coefficients and Jacobians describe that branch. The complete deployment rule includes transitions between raw, affine, and linear-program branches, each with its corresponding sensitivity.

\subsubsection{Active constraints in the final projection}

An explicit boundary in the same hypothetical model shows the change in sensitivity caused by non-negativity. Keep the total influent inventory equal to 10. The affine components obtained from Equation~\eqref{eq:icsor_small_back_projection} are
\begin{equation}
\equationname{Illustrative affine-projected components as functions of an operating input}
\label{eq:icsor_small_affine_branch}
c_{aff,1}=\frac{35}{6}+\frac{5a}{12},
\qquad
c_{aff,2}=\frac{25}{6}-\frac{5a}{12}.
\end{equation}
Their sum is 10. Near \(a=10\), the first component is positive, while the second reaches zero at \(a=10\). Immediately below this boundary the affine state is feasible and its derivative is \((5/12,-5/12)^\top\). Above the boundary the affine second component is negative. The positive-weight absolute projection returns \((10,0)^\top\), so its derivative with respect to \(a\) is \((0,0)^\top\). The final projected state therefore has a different derivative on each side of the boundary. The boundary is characterized by its two one-sided derivatives.

In the projected branch, the binding equalities are \(c_1+c_2=10\) and \(c_2=0\). Differentiating them gives \(\partial c_1/\partial a+\partial c_2/\partial a=0\) and \(\partial c_2/\partial a=0\). These two equations directly give the zero derivative. The general active-set calculation uses the same idea with all binding inventory and adjustment constraints.

The linear program can make a component concentration equal to zero or bind one of the absolute-adjustment inequalities. These binding constraints define an active set. If the optimal solution is unique and its nondegenerate optimal basis remains unchanged over a local neighborhood, the solution varies affinely with the relevant right-hand sides in that neighborhood. A branch-specific sensitivity can then be derived from that fixed basis.

At an active-set transition, the derivative can change abruptly. When several optimal projected states exist, the returned solution can also depend on the rule used to select among them. Local deployment sensitivity is assessed by differentiating a verified fixed basis where appropriate or by perturbing admissible inputs and resolving the entire deployment problem. The perturbation size and any branch changes are recorded as part of that sensitivity assessment.

This distinction prevents three different quantities from being conflated. A driver coefficient belongs to \(\widehat B\). A raw-state sensitivity passes through \(\widehat R^{-1}\). A deployed-state sensitivity also depends on the mass conservation reference and active projection constraints. Each quantity answers a different question about the model.

\subsubsection{Following sensitivity through projection}

Interpretation follows the complete prediction map. Coupling transforms the driver into component predictions. The reporting map combines components, and active output constraints change their local derivatives. Physical-unit interpretation identifies a driver coefficient, a raw prediction derivative, or a derivative of the final projected response \citep{Brunton2016,Shen2023}.

A second hypothetical calculation isolates the effect of projection using different raw equations. Let a dimensionless operating input be $a$, and let two hypothetical raw concentration equations be
\[
c_{raw,1}(a)=4+a,\qquad c_{raw,2}(a)=3+2a.
\]
The raw first-component slope is one and the second-component slope is two. At $a=2$, the concentrations are 6 and 7, with a total of 13. Suppose the required total is ten and the projection minimizes the sum of squared adjustments with equal weights. If the first adjustment is $d_1$ and the second is $d_2$, mass conservation requires $d_1+d_2=-3$. Substitute $d_2=-3-d_1$ into the squared-adjustment objective and differentiate,
\[
\begin{aligned}
d_1^2+d_2^2&=d_1^2+(-3-d_1)^2,\\
\frac{\mathrm d}{\mathrm d d_1}\bigl[d_1^2+(-3-d_1)^2\bigr]
&=4d_1+6=0.
\end{aligned}
\]
Both adjustments are $-1.5$, giving a state satisfying mass conservation with concentrations 4.5 and 5.5. For a general $a$, the excess total is $(4+a)+(3+2a)-10=3a-3$. Equal-weight projection subtracts one half of that excess from each coordinate,
\[
\begin{aligned}
c_{aff,1}(a)&=4+a-\frac{3a-3}{2}=5.5-0.5a,\\
c_{aff,2}(a)&=3+2a-\frac{3a-3}{2}=4.5+0.5a.
\end{aligned}
\]
The projected slopes are $-0.5$ and $0.5$. The first changes sign even though its raw slope is positive. The sum of the projected slopes is zero, consistent with a total that does not depend on $a$.

Non-negativity introduces another stage. The affine first component reaches zero at $a=11$. For $a>11$, the affine state lies outside the segment satisfying mass conservation and non-negativity. In this equal-weight two-component example, the nearest jointly feasible state is then $(0,10)^\top$. Both local output slopes are zero while that boundary remains active. The slope changes at the transition point. This calculation illustrates why a raw coefficient, an affine sensitivity, and a final constrained sensitivity have distinct meanings.

Figure~\ref{fig:assessment_branch_sensitivity} plots the hypothetical first-component equations and the slopes calculated from them. The left panel shows the concentration returned at each stage. The right panel shows the corresponding local rate of change with respect to the operating input.

\begin{figure}[htbp]
\centering
\includegraphics[width=\textwidth]{ch7_concept_branch_sensitivity.pdf}
\caption{Hypothetical first-component response and sensitivity before and after mass conservation and non-negativity enforcement. Open circles mark the two one-sided checked slopes at the active-boundary transition.}
\label{fig:assessment_branch_sensitivity}
\end{figure}

Before the boundary becomes active, the affine and checked curves coincide. Their negative first-component slope differs from the positive raw slope because mass conservation redistributes the operating response between the two components. After the boundary becomes active, the checked concentration stays at zero and its local slope is zero. The transition is characterized by the two one-sided checked slopes.

\subsection{Assessment design for the structured regression}

\subsubsection{Common inputs and comparator models}

The structured-regression assessment uses a design collection of 10,000 accepted steady-state reactor cases. Every method receives the same 22 physical inputs and predicts the same 20 effluent components. Its reported composites are obtained from the common composition matrix. The input domain, reaction model, simulation bounds, and acceptance criterion remain those of the declared reactor benchmark. The comparison therefore changes the predictor while holding the modeled process and prediction task fixed.

The nine unconstrained comparators are XGBoost, LightGBM, CatBoost, AdaBoost, random forest, support vector regression, \(k\)-nearest neighbors, partial least squares, and a fully connected multilayer perceptron. These methods represent boosted and randomized tree ensembles, kernel regression, local neighbor averaging, latent linear regression, and nonlinear neural regression. Their foundations include the ensemble construction of \citet{Breiman2001}, the boosting formulations of \citet{Drucker1997}, \citet{ChenGuestrin2016}, \citet{Ke2017}, and \citet{Prokhorenkova2018}, and the latent regression of \citet{Wold2001}. The comparators have no invariant penalty, auxiliary non-negative fitted-state matrix, or learned component-coupling equation during fitting. The latent components used by partial least squares are statistical input combinations. The shared 20-component prediction target remains the physical effluent state for every comparator.

The neural comparator has hidden layers with 128, 64, and 32 units and logistic activation. Dense connections map the 22 inputs through these layers to 20 component outputs. Its capacity is evaluated under the same samples and reported-composite map as the other predictors. A transparent polynomial and a neural network can therefore differ in their internal structure without changing the supervised target.

\subsubsection{Training partitions and hyperparameter selection}

The primary assessment uses a row-wise 80\% training and 20\% held-out split with random seed 42. The specified collection gives a training pool of 8000 rows and a held-out set of 2000 rows. Hyperparameter selection uses only the training pool. Half of that pool forms a 4000-row tuning subset. A further 80\% and 20\% split gives 3200 tuning-training rows and 800 tuning-validation rows. The tuning design uses one internal validation split.

Each predictor receives 100 hyperparameter trials under a seeded tree-structured Parzen estimator with seed 42. No trial pruning or wall-clock timeout is imposed. \citet{Bergstra2011} develop this search approach, while \citet{Akiba2019} provide the associated hyperparameter optimization framework. The selection criterion is mean squared error in the shared physical component space on the internal validation set. Selected settings are held fixed for the primary assessment and the repeated sample-size design. The final primary fit uses the full 8000-row training pool.

For the unconstrained comparators, input and output standardization parameters are estimated from the relevant training partition. Predictions are transformed back to physical component coordinates before any invariant check, reporting map, or error assessment. ICSOR remains in physical coordinates so that its driver, invariant operator, and deployed constraints use the same component convention. Internal-validation transformations are estimated from tuning-training rows only. The training partitions determine the means, standard deviations, and hyperparameters. Projection weights are fixed before held-out assessment.

Table~\ref{tab:icsor_comparator_settings} gives the specified numerical configurations. These values define fitted model settings rather than measured predictive performance. The internal search compares candidate configurations according to the stated validation objective.

\begin{table}[htbp]
\centering\small
\caption{Specified configurations of the nine unconstrained comparators.}\label{tab:icsor_comparator_settings}
\begin{tabular}{p{0.19\textwidth}p{0.74\textwidth}}
\toprule
Comparator & Configuration\\
\midrule
XGBoost & 309 trees, learning rate 0.125, maximum depth 3, minimum child weight 6.84, row fraction 0.915, feature fraction 0.828, absolute penalty \(1.18\times10^{-3}\), and squared penalty 0.382.\\
LightGBM & 302 trees, learning rate 0.0809, 16 leaves, maximum depth 8, minimum leaf count 14, row fraction 0.845, feature fraction 0.869, absolute penalty \(3.17\times10^{-3}\), and squared penalty \(3.36\times10^{-4}\).\\
CatBoost & 236 boosting iterations, learning rate 0.172, depth 4, leaf squared penalty 0.0266, bagging temperature 0.0277, and random strength 0.0657.\\
AdaBoost & 236 estimators, learning rate 0.999, and squared regression loss.\\
Random forest & 390 trees, maximum depth 14, minimum split count 3, minimum leaf count 1, feature fraction 0.500, and no bootstrap resampling.\\
Support vector regression & Polynomial kernel of degree 4, regularization parameter 0.118, insensitive-error width 0.0458, constant offset 0.889, and solver tolerance \(8.96\times10^{-4}\). The kernel scale is the reciprocal of the product of input dimension and training input variance.\\
\(k\)-nearest neighbors & 15 neighbors, inverse-distance weighting, Manhattan distance, and a k-dimensional search tree with leaf size 40.\\
Partial least squares & 6 latent components, maximum 887 fitting iterations, and tolerance \(9.91\times10^{-5}\).\\
Multilayer perceptron & Hidden widths \((128,64,32)\), logistic activation, adaptive moment estimation, squared penalty \(6.33\times10^{-4}\), batch size 64, initial learning rate \(1.08\times10^{-3}\), maximum 500 iterations, and tolerance \(1.63\times10^{-5}\). Training rows are not shuffled, and early stopping is disabled.\\
\bottomrule
\end{tabular}
\end{table}

\subsubsection{Scoring states and physical diagnostics}

The scored ICSOR output is the final checked component state. For comparator accuracy assessment, the raw component prediction receives the same affine projection for mass conservation in Equation~\eqref{eq:icsor_affine_solution}. The composition map is then applied to the projected state. This comparison describes the accuracy of a checked structured prediction and comparator predictions after affine projection. Comparator accuracy scoring uses the affine state, while ICSOR accuracy scoring uses the final state from checked deployment.

The numerical assessment used an operator \(A_{num}\) obtained by rounding the null-space basis, setting negligible coefficients to zero, and dividing each row by its first nonzero coefficient. The fitting, affine projection, and constrained deployment calculations used this operator in place of the ideal operator \(A\). The reported violation rates use these adopted numerical balances and the stated operator scaling.

Physical diagnostics use the direct raw states of the comparators. For ICSOR, diagnostics distinguish the raw coupled state, affine state, and final deployed state. This stage separation identifies the component state assessed by each diagnostic. The empirical balance diagnostic is \(v_{mc}^{num}(c)=\|A_{num}(c-c_{in})\|_2\). The non-negativity diagnostic remains \(v_{nn}\) in Equation~\eqref{eq:icsor_deployment_diagnostics}. At each stage, the violation frequency for diagnostic \(v\) is
\begin{equation}
\equationname{Percentage frequency of a physical-constraint violation}
\label{eq:icsor_violation_frequency}
f_v=\frac{100}{n}\sum_{i=1}^{n}
\mathbf{1}\{v(c_i)>10^{-10}\},
\end{equation}
where \(n\) is the number of assessed observations and \(v\) is either \(v_{mc}^{num}\) or \(v_{nn}\). Frequencies, residual magnitudes, adjustment distances, and deployment-branch use describe different aspects of the same prediction procedure. Predictive error is reported alongside these physical diagnostics.

\subsubsection{Composite error measures}

The differences among the five measures can first be computed for two hypothetical observations of one reported quantity. Take reference values \((2,4)\) and predictions \((3,2)\). Their residuals are \((1,-2)\). Squaring them gives \((1,4)\), while taking absolute values gives \((1,2)\). The mean squared error is \((1+4)/2=2.5\), its square root is approximately 1.5811, and the mean absolute error is \((1+2)/2=1.5\). These measures describe different summaries of the same two errors.

The absolute relative errors are \(1/2\) and \(2/4\), so their mean is 0.5. The reference mean is 3. Its total squared variation is \((2-3)^2+(4-3)^2=2\). The prediction residual squared sum is 5. Therefore the coefficient of determination is \(1-5/2=-1.5\). Its negative value says that the hypothetical predictions have greater squared error than assigning the reference mean to both observations. These values illustrate the definitions on a hypothetical pair of observations.

For reported quantity \(q\in\{1,2,3,4\}\), define the prediction residual \(e_{iq}=\widehat y_{iq}-y_{iq}\). The five measures are mean squared error, root mean squared error, mean absolute error, mean absolute relative error, and the coefficient of determination. They are computed per reported quantity and summarized across the four quantities. The squared and absolute measures are
\begin{equation}
\equationname{Per-composite squared, root, and absolute error measures}
\label{eq:icsor_composite_error_metrics}
E_{2,q}=\frac{1}{n}\sum_i e_{iq}^{2},
\qquad E_{root,q}=\sqrt{E_{2,q}},
\qquad E_{1,q}=\frac{1}{n}\sum_i|e_{iq}|.
\end{equation}
For strictly nonzero reference values, the relative measure is
\begin{equation}
\equationname{Mean relative absolute error of a reported composite}
\label{eq:icsor_relative_error}
E_{rel,q}=\frac{1}{n}\sum_i\frac{|e_{iq}|}{|y_{iq}|}.
\end{equation}
It is reported as a ratio. Multiplication by 100 converts it to a percentage. A zero reference value makes this expression undefined, and a very small value can dominate it. Any zero-reference convention is declared explicitly before comparing predictors. Targets with undefined relative error are identified under that convention.

The per-target coefficient of determination is
\begin{equation}
\equationname{Coefficient of determination for a reported composite}
\label{eq:icsor_coefficient_determination}
R_q^2=1-
\frac{\sum_i e_{iq}^{2}}
{\sum_i(y_{iq}-\overline y_q)^2}.
\end{equation}
A target with zero reference variance makes this expression undefined. A negative value is possible when the predictor has a greater squared error than the reference-mean predictor. This coefficient measures prediction error relative to the target variance.

The aggregate squared error pools the four composite residuals. Its root is
\begin{equation}
\equationname{Root mean squared error pooled across four reported composites}
\label{eq:icsor_aggregate_root_error}
E_{root}=\left(\frac{1}{4n}\sum_{q=1}^{4}\sum_{i=1}^{n}e_{iq}^{2}\right)^{1/2}.
\end{equation}
Aggregate absolute and relative errors use the corresponding mean over observations and composites. Aggregate \(R^2\) uses the mean of the four target-specific coefficients. The composite coordinates retain their own physical units. Their pooled error is a numerical benchmark summary under that coordinate convention. Per-target scores remain necessary to show which reported quantities contribute to an aggregate comparison.

Pooling before taking the square root also matters. If four hypothetical target mean squared errors are 1, 9, 1, and 9, the pooled mean squared error is 5 and its root is \(\sqrt5\), approximately 2.2361. The arithmetic mean of the four target root errors is instead \((1+3+1+3)/4=2\). These are different summaries. Equation~\eqref{eq:icsor_aggregate_root_error} specifies the first calculation.

The primary ordering measure is aggregate root mean squared error. The other measures provide complementary assessments. Squared error emphasizes larger residuals, absolute error gives average deviation, relative error emphasizes deviations in relation to the reference magnitude, and \(R^2\) compares the squared residual with reference variation. Predictor orderings are compared separately under each measure.

\subsubsection{Repeated sample-size assessment}

The sample-size design uses
\begin{equation}
\equationname{Observation counts for structured-regression sample-size assessment}
\label{eq:icsor_assessment_sizes}
n\in\{500,1450,2400,3350,4300,5250,6200,7150,8100,9050,10000\}.
\end{equation}
At each size, ten subsamples are drawn without replacement from the design collection. Deterministic seeds are derived from seed 42. Each run uses a row-wise 80\% training and 20\% held-out split. All ten predictors receive the same sampled rows and partitions. The design specifies 110 fits per predictor across the size schedule. The assessment uses repeated random splitting.

The fixed hyperparameters isolate sensitivity to available observations under one model specification. The mean and spread of each held-out score are summarized across the ten runs. Learning curves describe predictive performance across the entire sample-size schedule. \citet{Geman1992} explain the relation between flexibility and the bias-variance problem, and \citet{VieringLoog2023} review the varied shapes of learning curves.

The learning-curve area can first be calculated over a reduced hypothetical schedule. Suppose an error measure has values 6, 4, and 3 at sample sizes 100, 200, and 400. Between the first two sizes, a trapezoid has width 100 and mean height \((6+4)/2=5\), giving area 500. Between the next two, its width is 200 and mean height \((4+3)/2=3.5\), giving area 700. The total area is 1200. Dividing by the total width \(400-100=300\) gives a normalized area of 4. A larger interval contributes more to this average than a smaller interval, even when each has one pair of endpoints. This hypothetical schedule illustrates the area calculation for unequal sample-size intervals.

For successive retained sizes \(n_1,\ldots,n_S\), write the mean error of predictor \(m\) under measure \(j\) as \(\overline M_{m,j}(n_s)\). The general numerical calculation is
\begin{equation}
\equationname{Trapezoidal approximation of normalized learning-curve area}
\label{eq:icsor_learning_area_trapezoids}
\frac{1}{n_{max}-n_{min}}
\sum_{s=1}^{S-1}(n_{s+1}-n_s)
\frac{\overline M_{m,j}(n_s)+\overline M_{m,j}(n_{s+1})}{2}.
\end{equation}
Let \(\overline M_{m,j}(n)\) denote the mean held-out measure \(j\) for predictor \(m\) at sample size \(n\). The normalized area under its learning curve is
\begin{equation}
\equationname{Integral definition of normalized learning-curve area}
\label{eq:icsor_learning_area}
L_{m,j}=\frac{1}{n_{max}-n_{min}}
\int_{n_{min}}^{n_{max}}\overline M_{m,j}(n)\,\mathrm{d}n.
\end{equation}
The trapezoidal rule evaluates this integral over the stated sample sizes. Normalization by the sample-size interval gives an average height of the learning curve and preserves the units of measure \(j\). The normalization denominator is the sample-size interval. \citet{Cawley2011} and \citet{Guyon2011} provide related trajectory summaries in active learning. The present design varies the number of randomly sampled observations.

A rank discards the numerical distance between scores and keeps their order. For four hypothetical predictor errors 3.0, 3.0, 4.5, and 6.0, the dense ranks are 1, 1, 2, and 3. The tied first two predictors receive the same rank. There is no skipped rank after the tie. The difference between the third and fourth errors does not influence the difference between their ranks. If one predictor receives ranks 1, 2, and 2 in three illustrative comparisons, its mean rank is \((1+2+2)/3=5/3\). The actual assessment averages over the stated combinations of target, measure, and size.

For each measure \(j\), reported quantity \(q\), and retained size \(n_s\), predictors receive a dense rank \(r_{m,j,q}(n_s)\) from their mean held-out score. Lower squared, root, absolute, and relative errors receive better ranks. Larger \(R^2\) values receive better ranks. Equal scores receive equal ranks, and the next distinct score receives the next consecutive rank. With \(S=11\) sizes and five measures, the summaries are
\begin{equation}
\equationname{Mean dense ranks over quantities, measures, and sample sizes}
\label{eq:icsor_dense_rank}
\overline r_{m,j}=\frac{1}{4S}\sum_{q=1}^{4}\sum_{s=1}^{S}r_{m,j,q}(n_s),
\qquad
\overline r_m=\frac{1}{5}\sum_{j=1}^{5}\overline r_{m,j}.
\end{equation}
These summaries describe relative ordering across the stated comparisons. \citet{Demsar2006} develops broader rank-based model comparisons, while \citet{Shevchenko2024} emphasize variability in benchmarking. Mean ranks and absolute errors are reported together to describe both ordering and error magnitude.

At the largest sample size, the generalization gap is
\begin{equation}
\equationname{Held-out minus training root-error gap}
\label{eq:icsor_generalization_gap}
\Delta E_{root,m}=E_{root,m}^{held\text{-}out}-E_{root,m}^{training}.
\end{equation}
A positive gap indicates greater held-out error. A small gap can accompany either accurate prediction or high-error underfitting, so it is interpreted with the corresponding absolute scores. Training time is recorded for every repeated fit and summarized as a distribution across runs. Deployment effort is assessed separately through prediction time, projection-branch use, and constrained-solve effort. This separation makes the cost of learning distinct from the cost of returning a checked state.

For a hypothetical predictor with training error 0.5 and held-out error 0.8, the gap is 0.3. Another with errors 5.0 and 5.1 has a gap of only 0.1. The second has the smaller gap but much larger absolute errors. This arithmetic shows how absolute error and the training-to-held-out gap describe different aspects of prediction.

\subsection{Interpretability assessment}

The interpretability assessment examines the six driver blocks, the coupling matrix, and the corresponding component and composite raw-regression maps. It compares fitted response patterns with the declared process domain. The reported coefficient summaries describe the response structure of one fitted model. \citet{Brunton2016} connect interpretable nonlinear terms to the problem of recovering structure from data, while \citet{Machado2014} show how operational and influent uncertainty complicate parameter identification in activated sludge models.

The driver and coupling parameters can trade off in the relation \(Rc=B\phi\). Different joint stationary solutions can predict similar component states with different coefficients. Symmetry removes the redundancy of ordered self-quadratic pairs. Coefficient signs and magnitudes are interpreted together with the complete fitted response.

Prediction error and physical diagnostics are assessed separately for the raw regression, the affine state, and the final deployed state. The comparison distinguishes what the regression learns from what the projection enforces. The magnitude of the projection is also examined because a feasible deployed state can lie far from its raw prediction. Computational assessment records fitting effort and deployment effort separately. These measures describe the cost of fitting the structured regression and enforcing its physical conditions.

The assessment evaluates steady-state component prediction in the stated input domain and control volume. The raw regression supplies an inspectable approximation, and checked deployment enforces the adopted numerical invariants and non-negativity on the returned state.

\section{Results}
\label{sec:icsor_results}

\subsection{Accuracy on the fixed held-out partition}

Table~\ref{tab:icsor_fixed_results} reports the four-composite scores on the 2000 held-out cases. ICSOR used its final checked state. The other predictors used their affine-aligned states for accuracy scoring. The multilayer perceptron had the lowest pooled root mean squared error of 4.38. LightGBM followed at 5.30. Support vector regression and ICSOR were closely spaced at 5.97 and 5.98. These scores pool the native composite coordinates under Equation~\eqref{eq:icsor_aggregate_root_error}. The pooled scores retain the native composite-coordinate convention.

\begin{table}[htbp]
\centering\small
\caption{Fixed-split held-out composite performance. ICSOR is scored after checked deployment. Comparators are scored after affine alignment without a non-negativity constraint. Relative absolute error is a ratio rather than a percentage.}
\label{tab:icsor_fixed_results}
\begin{tabular}{p{4.7cm}rrrr}
\toprule
Predictor & $E_{root}$ & $E_1$ & Mean $R^2$ & $E_{rel}$\\
\midrule
Multilayer perceptron & 4.38 & 2.55 & 0.9915 & 0.0123\\
LightGBM & 5.30 & 3.21 & 0.9806 & 0.0203\\
Support vector regression & 5.97 & 3.49 & 0.9713 & 0.0244\\
ICSOR & 5.98 & 3.70 & 0.9700 & 0.0256\\
XGBoost & 6.12 & 3.81 & 0.9733 & 0.0244\\
CatBoost & 6.45 & 4.07 & 0.9727 & 0.0250\\
AdaBoost & 21.27 & 13.15 & 0.8643 & 0.0653\\
$k$-nearest neighbors & 23.61 & 13.87 & 0.8535 & 0.0557\\
Partial least squares & 28.24 & 17.05 & 0.7808 & 0.0681\\
Random Forest & 29.00 & 16.93 & 0.7934 & 0.0632\\
\bottomrule
\end{tabular}
\end{table}

ICSOR's pooled root error was 36.5\% higher than the multilayer perceptron's and 12.8\% higher than LightGBM's. It remained below XGBoost and CatBoost on this measure. The mean absolute relative error of 0.0256 corresponded to 2.56\%. The fixed-split comparison quantified the accuracy trade-off associated with the structured prediction procedure.

Figure~\ref{fig:icsor_fixed_results} compares root error, absolute error, and mean composite $R^2$. ICSOR ranked fourth by pooled root error on this partition.

\begin{figure}[htbp]
\centering
\includegraphics[width=\textwidth]{results/ch5/q01_fixed_split_accuracy.png}
\caption{Fixed-split composite accuracy on 2000 held-out cases. Teal bars identify ICSOR after checked deployment. Gray bars identify comparators after affine alignment. Pooled root and absolute errors retain the native composite-coordinate convention. The right panel averages four composite-specific $R^2$ values.}
\label{fig:icsor_fixed_results}
\end{figure}

Table~\ref{tab:icsor_target_results} separates the composite contributions. ICSOR's COD root error of 7.77 was 0.18 above LightGBM's and below that of the other non-neural predictors. Its TN and TSS errors were 3.55 and 8.36. The multilayer perceptron had the lowest COD, TN, and TSS errors. Every method had a TP error of 0.0361 after the specified alignment and reporting operations. TP therefore did not distinguish the predictors in this assessment.

\begin{table}[htbp]
\centering\small
\caption{Fixed-split root mean squared error for each composite. COD, TN, TP, and TSS use g COD m$^{-3}$, g N m$^{-3}$, g P m$^{-3}$, and g TSS m$^{-3}$, respectively.}
\label{tab:icsor_target_results}
\begin{tabular}{p{4.7cm}rrrr}
\toprule
Predictor & COD & TN & TP & TSS\\
\midrule
Multilayer perceptron & 6.48 & 1.58 & 0.0361 & 5.66\\
LightGBM & 7.59 & 2.77 & 0.0361 & 6.85\\
Support vector regression & 8.87 & 3.50 & 0.0361 & 7.20\\
ICSOR & 7.77 & 3.55 & 0.0361 & 8.36\\
XGBoost & 8.78 & 3.27 & 0.0361 & 7.88\\
CatBoost & 9.13 & 3.23 & 0.0361 & 8.53\\
AdaBoost & 38.30 & 5.58 & 0.0361 & 17.68\\
$k$-nearest neighbors & 38.24 & 4.23 & 0.0361 & 27.37\\
Partial least squares & 43.22 & 5.09 & 0.0361 & 36.00\\
Random Forest & 45.87 & 4.00 & 0.0361 & 35.26\\
\bottomrule
\end{tabular}
\end{table}

Figure~\ref{fig:icsor_composite_results} shows why the aggregate comparison needs separate target scores. COD and TSS contributed much larger numerical errors than TP. The separate target panels retain the constituent units used for each reported quantity.

\begin{figure}[htbp]
\centering
\includegraphics[width=\textwidth]{results/ch5/q02_composite_accuracy.png}
\caption{Fixed-split root mean squared error for COD, TN, TP, and TSS. Teal identifies ICSOR and gray identifies the comparators. Each panel has its own constituent unit and horizontal scale. All methods have the same reported TP error under the specified alignment and reporting procedure.}
\label{fig:icsor_composite_results}
\end{figure}

\subsection{Sample efficiency and generalization}

ICSOR had the lowest mean held-out root error at the smallest design size of 500 cases. Its value was 10.79, compared with 11.26 for XGBoost and 11.45 for LightGBM. The multilayer perceptron ranked fifth at 13.26. At the largest design size, the repeated-split means were 4.62 for the multilayer perceptron, 5.27 for LightGBM, 5.84 for support vector regression, and 5.94 for ICSOR. The ordering therefore changed as more training cases became available.

\begin{table}[htbp]
\centering\small
\caption{Repeated-split composite assessment over eleven design sizes and ten runs per size. Final error and the held-out minus training gap refer to $n=10{,}000$. The learning-curve area is normalized by the design-size interval. Lower error and curve-area values indicate better performance. The gap measures the separation between training and held-out error.}
\label{tab:icsor_scaling_results}
\begin{tabular}{p{4.3cm}rrrr}
\toprule
Predictor & \begin{tabular}{r}Final\\$E_{root}$\end{tabular} & \begin{tabular}{r}Normalized\\curve area\end{tabular} & Mean rank & $\Delta E_{root}$\\
\midrule
Multilayer perceptron & 4.62 & 6.75 & 2.5 & 0.65\\
LightGBM & 5.27 & 6.35 & 2.2 & 1.99\\
Support vector regression & 5.84 & 6.90 & 4.0 & 1.19\\
ICSOR & 5.94 & 6.33 & 4.3 & 0.22\\
XGBoost & 6.10 & 7.01 & 4.3 & 1.44\\
CatBoost & 6.26 & 7.60 & 5.0 & 0.69\\
AdaBoost & 21.42 & 20.29 & 8.2 & 0.23\\
$k$-nearest neighbors & 23.76 & 25.77 & 7.0 & 23.66\\
Partial least squares & 27.96 & 30.17 & 9.4 & 0.38\\
Random Forest & 28.78 & 30.31 & 8.2 & 18.71\\
\bottomrule
\end{tabular}
\end{table}

ICSOR had the lowest reported normalized root-error curve area of 6.33. LightGBM was close at 6.35 and had the best mean dense rank of 2.2. The multilayer perceptron had the second-best mean rank of 2.5. ICSOR and XGBoost each had a reported mean rank of 4.3. ICSOR led the curve-area comparison, while LightGBM led the mean-rank comparison.

ICSOR's held-out minus training root-error gap was 0.22. This was smaller than the gaps for the multilayer perceptron, LightGBM, XGBoost, and support vector regression. Partial least squares also had a relatively small gap of 0.38, but its held-out root error was 27.96. ICSOR combined a small gap with substantially lower absolute error than partial least squares.

Figure~\ref{fig:icsor_efficiency_results} presents the three trajectory summaries separately. ICSOR had the lowest root-error curve area, whereas LightGBM had the best mean rank. The two summaries answer different comparative questions.

\begin{figure}[htbp]
\centering
\includegraphics[width=\textwidth]{results/ch5/q03_sample_efficiency.png}
\caption{Normalized root-error curve area, mean dense rank, and held-out minus training root-error gap. Curve areas and ranks summarize eleven design sizes with ten repeated splits per size. The gap refers to the largest design size. Teal identifies ICSOR. The gap is reported alongside the held-out error.}
\label{fig:icsor_efficiency_results}
\end{figure}

Figure~\ref{fig:icsor_budget_results} compares the smallest and largest design sizes for four leading predictors. ICSOR led this group at 500 cases, while the multilayer perceptron led at 10,000 cases. The two endpoints show the change in the leading predictor between the smallest and largest designs.

\begin{figure}[htbp]
\centering
\includegraphics[width=\textwidth]{results/ch5/q04_data_budget_comparison.png}
\caption{Mean held-out pooled composite root error at design sizes 500 and 10,000 for ICSOR, XGBoost, LightGBM, and the multilayer perceptron. Values summarize ten repeated splits. The panels show repeated-split endpoint means at the two specified design sizes.}
\label{fig:icsor_budget_results}
\end{figure}

\subsection{Physical checks and deployment branches}

Table~\ref{tab:icsor_physical_results} distinguishes the final ICSOR output from the raw outputs of the comparators. The final structured predictions had no violations of the adopted numerical balances or non-negativity tolerance. Every comparator raw state violated the adopted balance tolerance. AdaBoost, Random Forest, and the nearest-neighbor model nevertheless had no negative raw states. The other raw non-negativity violation rates ranged from 46.00\% to 70.10\%.

\begin{table}[htbp]
\centering\small
\caption{Component-state violation frequencies on the 2000 held-out cases. ICSOR uses its final checked state. Comparators use their raw state, not the affine-aligned state used for accuracy scoring. Balance violations refer to $A_{num}$ and the declared $10^{-10}$ tolerance.}
\label{tab:icsor_physical_results}
\begin{tabular}{p{5.2cm}rr}
\toprule
Predictor & Balance violations (\%) & Negative states (\%)\\
\midrule
Multilayer perceptron & 100.0 & 58.20\\
LightGBM & 100.0 & 46.00\\
Support vector regression & 100.0 & 68.65\\
ICSOR & 0.0 & 0.00\\
XGBoost & 100.0 & 64.95\\
CatBoost & 100.0 & 70.10\\
AdaBoost & 100.0 & 0.00\\
$k$-nearest neighbors & 100.0 & 0.00\\
Partial least squares & 100.0 & 48.70\\
Random Forest & 100.0 & 0.00\\
\bottomrule
\end{tabular}
\end{table}

No raw coupled ICSOR prediction passed both physical checks. All failed the balance tolerance, although 21.7\% were non-negative. Affine projection restored the adopted balances for every case and also returned a non-negative state for 436 cases, or 21.8\%. Those cases were returned after the affine stage. The remaining 1564 cases, or 78.2\%, required the weighted linear projection. The solve used a mean of 23.5 iterations and a maximum of 40. Every final state passed both checks.

Figure~\ref{fig:icsor_deployment_results} keeps the comparator diagnostics and ICSOR's deployment stages separate. The rise from zero raw balance compliance to complete final compliance occurred through projection. The branch counts show that the constrained solve was used for most deployed predictions.

\begin{figure}[htbp]
\centering
\includegraphics[width=\textwidth]{results/ch5/q05_physical_deployment.png}
\caption{Physical diagnostics and ICSOR deployment on 2000 held-out cases. The left panel compares final ICSOR states with raw comparator states. The middle panel separates raw coupled, affine, and final ICSOR compliance. The right panel gives the numbers returned after affine projection and after constrained linear projection. Balance checks use the adopted numerical operator $A_{num}$.}
\label{fig:icsor_deployment_results}
\end{figure}

These stage counts identify the source of the final compliance. The affine stage restored the adopted balances, and the constrained projection supplied non-negativity for the remaining cases. The complete deployment procedure returned zero final violation rates.

\subsection{Computational effort}

At the smallest design size, mean ICSOR fitting time was 21.5 s. LightGBM, XGBoost, and the multilayer perceptron required 4.72, 3.16, and 4.99 s, respectively. At the largest size, ICSOR required 94.4 s. The faster tree ensembles required 4.13--8.13 s, while partial least squares and the nearest-neighbor model required 0.228 and 2.70 s. ICSOR was nevertheless faster to fit than AdaBoost at 117.7 s, support vector regression at 136.8 s, and the multilayer perceptron at 161.6 s.

The recursive constrained fitting procedure was more costly than the fastest comparators and faster than three comparators at the largest design size. Fitting time was assessed separately from deployment effort. The 78.2\% linear-projection activation rate showed that most fixed-split predictions required an optimization step after regression evaluation. The reported iteration counts described this deployment workload.

Figure~\ref{fig:icsor_fitting_results} separates the two design sizes on logarithmic time axes. The smaller-size panel compares four predictors, and the larger-size panel compares ICSOR with five alternatives. The comparison reports the time required for the repeated model fits.

\begin{figure}[htbp]
\centering
\includegraphics[width=\textwidth]{results/ch5/q06_fitting_effort.png}
\caption{Mean fitting times for the specified predictors at design sizes 500 and 10,000. Teal identifies ICSOR. Each panel uses a logarithmic seconds axis. The values concern repeated model fits and exclude deployment calculations.}
\label{fig:icsor_fitting_results}
\end{figure}

\subsection{Inspectable coefficient structure}

Table~\ref{tab:icsor_coefficient_results} reports the COD raw-composite coefficient counts and the shared coupling entries. The polynomial summary counted 276 unique candidates. Adding the 380 off-diagonal coupling candidates gave 656 entries in this display. The display combined the COD polynomial row with the shared coupling entries. The display retained 469 entries, or 71.5\%, under the specified post-fitting thresholds.

\begin{table}[htbp]
\centering\small
\caption{Coefficients exceeding the reporting thresholds in one fitted model. Polynomial blocks describe the COD row of $B_{\mathrm{raw,comp}}$. The coupling matrix is shared across component outputs and uses its own threshold. Symmetric quadratic blocks count unique upper-triangular terms.}
\label{tab:icsor_coefficient_results}
\begin{tabular}{p{6.2cm}rrr}
\toprule
Coefficient block & Candidates & Retained & Retained (\%)\\
\midrule
Baseline $b$ & 1 & 1 & 100.0\\
Operating terms $W_u$ & 2 & 2 & 100.0\\
Influent terms $W_{in}$ & 20 & 19 & 95.0\\
Operating curvature $\Theta_{uu}$ & 3 & 3 & 100.0\\
Operating-load interactions $\Theta_{uc}$ & 40 & 29 & 72.5\\
Influent curvature $\Theta_{cc}$ & 210 & 44 & 21.0\\
Shared off-diagonal coupling $\Gamma$ & 380 & 371 & 97.6\\
Combined display & 656 & 469 & 71.5\\
\bottomrule
\end{tabular}
\end{table}

The influent curvature block retained 44 of 210 unique terms. In contrast, 371 of 380 coupling entries exceeded their reporting threshold. The fitted structure therefore combined selective influent curvature with dense component coupling. Selectivity was concentrated in the influent-curvature block, while the coupling matrix was dense.

The largest absolute COD influent-curvature entry was the dissolved-oxygen self term of $-5.8097$. Each mirrored dissolved-oxygen and dinitrogen cell had value $-1.7411$. Their combined unordered interaction coefficient was $-3.4822$. Other prominent terms linked dissolved oxygen with metal phosphate, ammonia-oxidizing biomass, soluble phosphate, nitrite-oxidizing biomass, and nitrite. These values described the raw composite map after coupling and reporting.

Figure~\ref{fig:icsor_structure_results} displays the retained counts and selected oxygen-related curvature terms. Its interaction values combine the two mirrored cells. The comparison describes the fitted polynomial under its native input units.

\begin{figure}[htbp]
\centering
\includegraphics[width=\textwidth]{results/ch5/q07_coefficient_structure.png}
\caption{Coefficient structure of one fitted ICSOR model. The left panel separates entries above and below the display thresholds for the COD polynomial blocks and shared coupling. The right panel shows selected COD influent-curvature coefficients with mirrored interaction entries combined. Coefficient units depend on the corresponding input terms. The displayed values describe the raw composite mapping.}
\label{fig:icsor_structure_results}
\end{figure}

\section{Discussion}
\label{sec:icsor_discussion}

\subsection{A structured approximation under limited data}

ICSOR's lowest small-design error and lowest reported root-error curve area supported a data-efficiency benefit in the evaluated domain. Its explicit second-order terms defined a structured response approximation. The invariant-aware fitting objective and learned component coupling supplied additional structure. The combined design performed well when fewer simulated states were available. This interpretation is consistent with the role of model structure in hybrid learning discussed by \citet{Shen2023} and with the changing model orderings described by \citet{VieringLoog2023}.

The leading predictor depended on the assessment measure. LightGBM had the best mean rank across targets, metrics, and sizes. The multilayer perceptron had the lowest full-size error. ICSOR's curve area was 6.33, close to LightGBM's 6.35. Selection based on full-size composite error favored the multilayer perceptron under this protocol. An inspectable equation, data efficiency, and checked final states provided additional engineering criteria for selecting ICSOR.

The training-to-held-out gap and absolute error described complementary properties. Partial least squares combined a small gap with high prediction error. ICSOR combined a gap of 0.22 with substantially lower held-out error, while the leading flexible predictors achieved still lower absolute error. Reporting both measures clarified the relation between generalization and accuracy.

\subsection{Fitting guidance and final-state enforcement}

The physical assessment showed how the complete procedure supplied its returned state. The soft invariant penalty guided the fitted training states, and the learned coupled equation generated the raw prediction for each new input. Every fixed-split raw prediction required affine projection. The affine state passed both checks in 21.8\% of cases, and the weighted linear projection handled the remaining 78.2\%.

The complete deployment procedure returned zero final balance and non-negativity violations. Stage-wise diagnostics identified the affine and constrained projections as the source of this compliance. This separation between constrained fitting and output enforcement follows the approaches examined by \citet{Hansen2024} and \citet{KircherVotsmeier2025}.

The numerical assessment used the rounded and normalized operator $A_{num}$ with a tolerance of $10^{-10}$. Accuracy scoring used checked ICSOR states and affine-aligned comparator states. Physical comparator diagnostics used their raw states. These specified procedures produced the 36.5\% root-error difference between ICSOR and the multilayer perceptron. The shared alignment and reporting convention also produced the common TP error of 0.0361.

\subsection{What the coefficient structure explains}

The named coefficient blocks exposed operating dependence, loading dependence, curvature, and component coupling for direct inspection. This form of interpretability follows the preference for transparent models discussed by \citet{Rudin2019}. Coupling and the reporting map translated the component-driver coefficients into composite response equations.

The retained counts showed where the fitted representation was selective and where it remained dense. Most unique influent-curvature terms were below the display threshold, while almost all coupling entries exceeded their separate threshold. The display retained 469 of 656 entries. The complete fitted mapping retained its coefficients, with reporting thresholds controlling the detail visible in the display.

The prominent oxygen-related curvature terms were consistent with oxygen-sensitive reaction pathways in the adopted process model. \citet{StenstromSong1991} provide process context for oxygen dependence in nitrification, while \citet{Wentzel1992} describe interactions among nutrient-removal pathways. The dissolved-oxygen self coefficient and the oxygen-dinitrogen interaction therefore identified fitted response patterns that could be compared with the reaction model. Their contributions were interpreted in the native input units and through the complete coupled mapping.

The polynomial supplied analytical raw sensitivities, while the affine projection modified them to preserve the influent inventories. Equations~\eqref{eq:icsor_raw_jacobians} and \eqref{eq:icsor_affine_jacobians} described these two stages. The frequent constrained branch made active-set sensitivity relevant to the final returned response. Figure~\ref{fig:assessment_branch_sensitivity} demonstrated this change through an explicit example in which the first-component slope changed from positive to negative and then became zero at an active non-negativity boundary.

\subsection{Engineering implications}

The structured model combined competitive composite prediction, favorable low-data behavior, and final states that passed the imposed numerical checks. Its second-order terms and coupling map made the fitted response available for inspection. Checked deployment then enforced the adopted material accounting and component non-negativity before reporting COD, TN, TP, and TSS.

These properties involved distinct computational costs. ICSOR required more fitting time than the fastest tree and linear comparators. At 10,000 cases, it fitted faster than AdaBoost, support vector regression, and the multilayer perceptron. The constrained deployment solve was used for most fixed-split predictions. Separate fitting and deployment measurements therefore described the effort required to construct the approximation and return a checked state for engineering analysis.
